\section*{Chapitre \ref{ch:iv3:cpp} --- C++}
\markboth{\protect\omTocMark{}{Réponses~: C++}}{\protect\omTocMark{}{Réponses~: C++}}

\begin{solution}{iq:iv3:cpp:1}
\texttt{p = \&b} et \texttt{*q = 3} compilent, et les quatre autres ne compilent pas. Il faut lire une déclaration à partir du nom, vers l'extérieur et vers la gauche~: \texttt{p} est un pointeur vers un \texttt{const int}, et \texttt{q} est un pointeur \texttt{const} vers un \texttt{int}.

Ainsi, \texttt{p} peut pointer vers un autre objet et ne peut pas écrire par son intermédiaire. \texttt{q} peut écrire sans pouvoir changer de cible, et \texttt{r} ne peut faire ni l'un ni l'autre. \texttt{a} lui-même n'est pas constant, donc une lecture à travers \texttt{p} peut toujours voir une modification faite par \texttt{a} ou \texttt{q}~: \texttt{const int*} garantit seulement que ce chemin n'écrira pas.

\texttt{const int*} et \texttt{int const*} sont le même type, car la position de \texttt{const} par rapport à \texttt{int} n'a pas d'importance. Seule compte sa position par rapport à \texttt{*}.

\iqwrong{«~\texttt{const int* p} signifie que la valeur ne peut pas changer.~» Elle empêche seulement les écritures à travers \texttt{p}.}
\iqfollow{Et dans la signature d'une fonction~?}{\texttt{void f(int* const p)} a le même type que \texttt{void f(int* p)}~: le \texttt{const} de premier niveau porte sur la copie que la fonction possède de l'argument, donc il ne modifie pas la vue de l'appelant.}
\end{solution}

\begin{solution}{iq:iv3:cpp:2}
Il affiche \texttt{+global +static +auto -auto +auto -auto +heap end -static -global}. L'objet du tas n'est jamais détruit, puisque rien n'appelle \texttt{delete}, donc \texttt{-heap} n'apparaît jamais.

\texttt{global} a une durée de stockage statique et est construit avant \texttt{main}. \texttt{s} a aussi une durée de stockage statique, mais il est construit la première fois que l'exécution passe par sa déclaration, ce qui explique que \texttt{+static} n'apparaisse qu'une fois. \texttt{a} a un stockage automatique, il est construit et détruit à chaque appel, et il vit sur la pile. \texttt{*h} a un stockage dynamique sur le tas, et le pointeur \texttt{h} lui-même est automatique.

À la sortie, les objets statiques sont détruits dans l'ordre inverse de leur construction, donc \texttt{s} précède \texttt{global}. Les deux objets statiques résident dans le segment de données du programme, et non sur la pile ni sur le tas.

\iqwrong{«~Le système d'exploitation libère l'objet du tas à la sortie, donc son destructeur s'exécute.~» La mémoire est récupérée, mais aucun destructeur ne s'exécute. Un destructeur qui viderait un tampon de fichier ne le viderait jamais.}
\iqfollow{Le local statique est-il sûr pour les threads~?}{Depuis C++11, son initialisation l'est~: un second thread attend qu'elle soit terminée.}
\end{solution}

\begin{solution}{iq:iv3:cpp:3}
Seul l'accès par défaut diffère~: les membres et les bases d'une \texttt{struct} sont publics par défaut, et ceux d'une \texttt{class} sont privés. \texttt{S\{\}.x}, \texttt{C\{\}.get()} et \texttt{FromC\{\}.get()} compilent, tandis que \texttt{C\{\}.x} et \texttt{FromS\{\}.x} ne compilent pas.

\texttt{FromC} est une \texttt{struct}, donc il hérite publiquement, et \texttt{get} reste public. \texttt{FromS} est une \texttt{class}, donc il hérite en privé, ce qui masque \texttt{x} alors que \texttt{x} est public dans \texttt{S}. Le mode par défaut vient de la classe dérivée, et non de la base.

Les deux mots-clés peuvent avoir des méthodes, des constructeurs, des fonctions virtuelles et des modèles, et peuvent être mélangés librement dans une même hiérarchie. La disposition en mémoire ne dépend pas du mot-clé.

La convention courante est d'utiliser \texttt{struct} pour un bundle de données publiques sans invariant, et \texttt{class} lorsque les membres sont gardés privés afin de protéger un invariant.

\iqwrong{«~Une \texttt{struct} n'est que de simples données, et elle vit sur la pile.~» C'est la distinction de C\#, non celle de C++. L'endroit où vit un objet dépend de la manière dont il est créé.}
\end{solution}

\begin{solution}{iq:iv3:cpp:4}
Il affiche \texttt{2 3 2} et \texttt{2 4 4}. \texttt{r = b} affecte la valeur de \texttt{b} à \texttt{a}, car \texttt{r} n'est qu'un autre nom pour \texttt{a}, et une référence ne peut pas être redirigée. \texttt{p = \&b} fait pointer le pointeur ailleurs, et laisse \texttt{a} intact.

Ainsi, après \texttt{r = b}, \texttt{a} vaut 2, et l'affectation suivante \texttt{b = 3} ne le touche pas. Ensuite, \texttt{*p = 4} écrit dans \texttt{b}.

Deux autres différences. Une référence doit être liée au moment de sa déclaration, et elle n'est jamais censée être nulle, si bien qu'une fonction qui en reçoit une n'a pas besoin de vérifier. Un pointeur est en revanche un objet à part entière, dont on peut prendre l'adresse, qu'on peut ranger dans un conteneur et sur lequel on peut faire de l'arithmétique, alors que \texttt{\&r} donne simplement l'adresse de \texttt{a}.

\iqwrong{«~Une référence est plus sûre, car elle ne peut pas être pendante.~» Elle peut l'être. Une référence à une variable locale renvoyée par une fonction est pendante, tout comme un pointeur.}
\iqfollow{Quand un pointeur est-il le bon paramètre~?}{Lorsque «~rien~» est un argument valide, ou lorsque la fonction le conserve et qu'il peut être redirigé plus tard.}
\end{solution}

\begin{solution}{iq:iv3:cpp:5}
Il affiche 40, 8 et 10. Dans \texttt{main}, \texttt{a} est un tableau de 10 \texttt{int}, et \texttt{sizeof} donne 40 octets. En tant que paramètre, \texttt{int a[10]} signifie \texttt{int* a}, de sorte que la fonction voit un pointeur, de 8 octets, et le 10 entre crochets est ignoré.

Cette conversion implicite est héritée du C~: un tableau passé à une fonction se transforme en pointeur vers son premier élément. La longueur n'est pas transmise du tout, si bien que l'appelant pourrait fournir un tableau de 3, sans que le compilateur dise quoi que ce soit. g++ signale par défaut ce \texttt{sizeof}, raison pour laquelle le listing désactive cet avertissement sur une ligne non montrée.

Pour conserver la longueur, on prend un \texttt{std::span<int>} en C++20, qui porte le pointeur et la taille et accepte indifféremment tableaux et vecteurs. Ou bien on prend une référence au tableau, \texttt{int (\&a)[10]}, ou on utilise un modèle paramétré par la longueur. \texttt{std::size} fonctionne sur un vrai tableau parce que son type porte la longueur.

\iqwrong{«~\texttt{sizeof(a) / sizeof(a[0])} donne la longueur.~» Dans la fonction, il donne $8 / 4 = 2$, quoi qu'on ait passé.}
\iqfollow{Pourquoi \texttt{std::array} évite-t-il cela~?}{C'est une struct qui contient le tableau, si bien qu'on le passe par valeur ou par référence, sa taille figurant dans le type.}
\end{solution}

\begin{solution}{iq:iv3:cpp:6}
Il affiche \texttt{300 4}, puis \texttt{44}, puis \texttt{A 65 65}. \texttt{a + b} est de type \texttt{int}, et rien ici n'a de comportement indéfini.

Les types plus petits que \texttt{int} sont promus en \texttt{int} avant le calcul, de sorte que \texttt{a + b} vaut l'entier 300, et \texttt{auto} fait de \texttt{c} un \texttt{int} de 4 octets. Stocker 300 dans \texttt{uint8\_t} le réduit modulo 256, ce qui donne 44. Cette conversion vers un type non signé est bien définie. \texttt{-Wall -Wextra} n'émet aucun avertissement, et \texttt{-Wconversion} seul non plus avec g++ 11.4~: il faut \texttt{-Wconversion -Warith-conversion}.

Ici, \texttt{std::uint8\_t} est un alias de \texttt{unsigned char}, donc \texttt{cout} l'affiche comme un caractère, et 65 apparaît comme \texttt{A}. \texttt{int(e)} ou le plus unaire, qui effectue la promotion, affiche le nombre.

Le piège, dans les analyseurs de feed, c'est qu'un champ d'un octet journalisé avec \texttt{cout} s'affiche comme un caractère de contrôle, ou pas du tout.

\iqwrong{«~\texttt{a + b} déborde, puisque les deux font 8 bits.~» L'addition se fait dans \texttt{int}, donc il n'y a pas de débordement.}
\iqfollow{La promotion peut-elle provoquer un comportement indéfini~?}{Oui~: deux valeurs \texttt{uint16\_t} de 65535 se multiplient comme des \texttt{int} signés, et cela déborde.}
\end{solution}

\begin{solution}{iq:iv3:cpp:7}
Il affiche \texttt{3 2 3 1} et \texttt{1 3 0 3}. Les parenthèses choisissent un constructeur comme d'habitude, donc \texttt{a(3, 1)} contient trois copies de 1, et \texttt{c(3)} contient trois zéros. Les accolades privilégient un constructeur \texttt{initializer\_list} dès qu'il s'applique, donc \texttt{b\{3, 1\}} contient les deux éléments 3 et 1, et \texttt{d\{3\}} contient un seul élément.

\begin{sloppypar}La règle est que l'initialisation par accolades examine d'abord les constructeurs \texttt{initializer\_list}, et ne considère les autres que si aucun n'est viable. Ainsi, pour tout conteneur de nombres, les accolades signifient «~ces éléments~».\end{sloppypar}

Avec \texttt{std::vector<std::string> v\{3, \textquotedbl{}x\textquotedbl{}\}}, la liste ne peut pas se convertir, si bien que c'est le constructeur prenant un nombre et une valeur qui est utilisé, et l'on obtient trois chaînes de caractères. Les mêmes accolades changent donc de sens selon le type des éléments, ce qui compte dans les modèles.

\iqwrong{«~Les accolades sont la manière moderne, donc il faut les utiliser partout.~» Elles interdisent aussi les conversions restrictives, ce qui est bien, mais sur un conteneur de nombres elles font autre chose que les parenthèses.}
\iqfollow{Quelle bonne habitude recommandez-vous~?}{Utiliser des parenthèses pour une taille et des accolades pour une liste d'éléments, et ne pas les mélanger dans du code générique sans réfléchir.}
\end{solution}

\begin{solution}{iq:iv3:cpp:8}
Il affiche \texttt{+a +b body -b -a}. Les membres sont initialisés dans l'ordre de leur déclaration dans la classe, quel que soit l'ordre de la liste d'initialisation, et détruits dans l'ordre inverse. \texttt{-Wreorder}, que \texttt{-Wall} inclut, signale la discordance.

Le listing désactive cet avertissement sur une ligne non montrée, car sinon il ferait échouer une construction avec \texttt{-Werror}.

Cela compte lorsqu'un initialiseur dépend d'un autre. Avec \texttt{Pair() : b(a.name), a(\textquotedbl{}a\textquotedbl{})}, \texttt{b} est construit en second, si bien que celui-ci fonctionne par hasard. La dépendance inverse est le bogue~: si l'initialiseur de \texttt{a} lit \texttt{b}, il lit un objet qui n'existe pas encore, ce qui est un comportement indéfini, alors que la lecture de la liste laisse croire le contraire.

La raison d'un ordre fixe tient à la destruction. Il peut y avoir plusieurs constructeurs avec des listes différentes, mais un seul destructeur, et celui-ci doit défaire la construction en sens inverse.

\iqwrong{«~C'est la liste qui décide, puisqu'elle s'exécute de gauche à droite.~» Les classes de base viennent d'abord, puis les membres dans l'ordre de déclaration.}
\iqfollow{Et le corps du constructeur~?}{Il s'exécute une fois que tous les membres sont construits, si bien qu'une affectation dans le corps est une seconde écriture, et non une initialisation.}
\end{solution}

\begin{solution}{iq:iv3:cpp:9}
Il affiche \texttt{int}, \texttt{pointer}, \texttt{pointer}, \texttt{pointer}, \texttt{int}. \texttt{nullptr} est de type \texttt{std::nullptr\_t}, qui se convertit en tout type pointeur et non en \texttt{int}. Le littéral 0 est un \texttt{int}, donc il sélectionne la surcharge \texttt{int}.

Avant C++11, la valeur nulle s'écrivait 0 ou \texttt{NULL}, et l'un ou l'autre pouvait choisir silencieusement une surcharge pour les entiers. Un littéral 0 cesse aussi d'être un pointeur nul une fois qu'il passe par un modèle~: dans \texttt{relay(0)}, \texttt{T} vaut \texttt{int}, et une variable \texttt{int} contenant 0 ne peut pas se convertir en pointeur, si bien que, s'il n'existait que la surcharge pour pointeur, \texttt{relay(0)} ne compilerait pas. \texttt{relay(nullptr)} conserve le type, et celui-ci se convertit toujours.

Ainsi, \texttt{nullptr} exprime «~aucun objet~» dans le système de types, et il se conserve lorsqu'on le transmet. Sa valeur n'a pas d'importance~: sous Linux x86-64, le pointeur nul est l'adresse dont tous les bits sont à zéro, mais aucun code ne devrait en dépendre.

\iqwrong{«~0 et \texttt{nullptr} sont la même chose.~» Dans \texttt{f(0)}, ils appelleraient des fonctions différentes, ce qui est tout l'intérêt.}
\end{solution}

\begin{solution}{iq:iv3:cpp:10}
Il affiche \texttt{5: 1 2 3 0 3} puis \texttt{3}. \texttt{std::remove} travaille sur des itérateurs, donc il ne peut que déplacer les éléments. Il ne peut pas changer la taille d'un conteneur, puisqu'il ne voit jamais le conteneur. En C++20, \texttt{std::erase(v, 0)} effectue les deux étapes.

Les éléments conservés sont décalés vers l'avant, et la fonction renvoie la nouvelle fin logique. Ce qui se trouve au-delà est valide mais non spécifié, ce qui explique que la queue lue ici soit \texttt{0 3}~: ce sont les anciennes positions, laissées telles quelles. Pour un type disposant d'un déplacement, elles peuvent se trouver dans un état déplacé.

La forme classique est \texttt{v.erase(std::remove(v.begin(), v.end(), 0), v.end())}, qui est linéaire, chaque élément conservé n'étant déplacé au plus qu'une fois. C'est bien meilleur que d'effacer un élément à la fois, ce qui est quadratique.

\iqwrong{Utiliser \texttt{v.end()} après \texttt{std::remove} comme si l'itérateur avait été déplacé. Ce n'est pas le cas, si bien qu'une boucle allant jusqu'à \texttt{v.end()} voit les restes.}
\iqfollow{Et pour une \texttt{std::list}~?}{Elle possède sa propre méthode membre \texttt{remove}, qui détache les nœuds et change bien la taille.}
\end{solution}

\begin{solution}{iq:iv3:cpp:11}
Il affiche 0 puis 1.8. Dans \texttt{accumulate}, le type du cumul est celui de la valeur initiale~; un littéral 0 en fait donc un \texttt{int}, et chaque étape tronque.

Les étapes valent donc $0 + 0.6 = 0.6$, stockées dans un \texttt{int} comme 0, trois fois de suite. Le résultat est 0, et même pas 1 ou 2, puisque rien ne se reporte jamais. Avec des exécutions de 0.6, le total reste à 0 pour toujours, ce qui serait grave sur une position.

\texttt{-Wall -Wextra} ne signalent rien, puisque la conversion a lieu dans un modèle de bibliothèque, dans un en-tête système, où les avertissements sont désactivés par défaut. Écrire le départ sous la forme \texttt{0.0}, utiliser \texttt{std::reduce} avec un départ typé, ou fixer soi-même le type~: \texttt{accumulate(b, e, double\{\})}.

Le même piège existe avec un départ à \texttt{0} sur une somme de valeurs \texttt{long}, où le résultat de chaque étape, en \texttt{long}, est ramené à un \texttt{int}, si bien qu'un total important sort silencieusement faux.

\iqwrong{«~Cela arrondit, donc la réponse est proche de 2.~» Chaque étape tronque vers zéro et repart de ce qui reste.}
\end{solution}

\begin{solution}{iq:iv3:cpp:12}
Il affiche \texttt{1 2}, \texttt{10 20} et \texttt{10 7}. La première boucle modifie une copie, et la seconde modifie les éléments. \texttt{auto} omet la référence et le \texttt{const} de premier niveau, si bien que \texttt{a} est un simple \texttt{int}, et y écrire laisse \texttt{v} intact.

Pour les conserver, il faut le dire~: \texttt{auto\& b = r} déduit \texttt{const int\&}, et \texttt{const auto\&} l'écrit explicitement, les deux se liant à l'élément. \texttt{decltype(auto)} conserve le type exact de l'initialisateur.

Pour une boucle basée sur un intervalle, utiliser \texttt{const auto\&} pour tout ce qui peut coûter cher à copier, comme une chaîne ou une grande structure, et \texttt{auto\&} pour modifier les éléments. Le simple \texttt{auto} convient aux petites valeurs, comme un entier ou un \texttt{string\_view}, dont la copie est peu coûteuse et où le code ne voit pas d'aliasing.

\iqwrong{\texttt{for (auto x : orders) x.qty = 0;} «~Ça compile, donc ça marche.~» Cela met à zéro des copies, sans aucun avertissement, et les ordres restent inchangés.}
\iqfollow{Pourquoi \texttt{const auto\&} peut-il poser problème sur une map~?}{Avec un type d'élément qui ne correspond pas, comme \texttt{const std::pair<int, double>\&} pour le type \texttt{pair<const int, double>} de la map, chaque élément est converti en temporaire puis lié~: une copie non remarquée à chaque étape.}
\end{solution}

\begin{solution}{iq:iv3:cpp:13}
\texttt{b = a} copie le pointeur, si bien que les deux objets possèdent un seul tableau. Quand \texttt{b} est détruit, le tableau est libéré, et AddressSanitizer signale un \texttt{heap-use-after-free} à \texttt{a.data[0] = 2.0}. Si cette ligne était absente, le destructeur de \texttt{a} libérerait le tableau une seconde fois.

Le compilateur écrit lui-même les opérations de copie, sauf indication contraire, et il copie les membres un à un, ce qui, pour un pointeur brut, signifie l'adresse et non la mémoire pointée.

La règle des cinq~: si une classe a besoin d'un destructeur, d'un constructeur de copie, d'un opérateur d'affectation par copie, d'un constructeur de déplacement ou d'un opérateur d'affectation par déplacement personnalisé, elle doit probablement envisager les cinq. Le \Cref{lst:iv3:cp:bufferfive} le fait. L'affectation par valeur sert à la fois à copier et à déplacer, et son échange rend l'auto-affectation sans risque. \texttt{noexcept} sur les déplacements permet aux conteneurs de les employer.
\omcode{code/interviews-3/03-cpp/cpp/buffer_five.cpp}{1}{24}{Les cinq fonctions, avec affectation par valeur.}\label{lst:iv3:cp:bufferfive}

Mieux encore, la règle du zéro~: détenir un \texttt{std::vector<double>} et n'en écrire aucune. Les versions du compilateur sont alors correctes, puisque chaque membre se gère lui-même. Écrire les cinq est réservé aux rares classes dont le seul rôle est de posséder une ressource.

\iqwrong{«~Déclarer le constructeur de copie privé et c'est réglé.~» Cela empêche la copie, et en C++ moderne on écrirait \texttt{= delete}~; c'est acceptable pour un type qui ne doit pas être copié, mais il faut toujours décider des déplacements.}
\end{solution}

\begin{solution}{iq:iv3:cpp:14}
Il affiche \texttt{0 1} puis \texttt{1 0}. Un destructeur déclaré par l'utilisateur fait que le compilateur ne génère ni le constructeur de déplacement ni l'affectation par déplacement, si bien que \texttt{std::move} se rabat sur le constructeur de copie. Le programme reste correct, et il n'a perdu que de la vitesse. La correction consiste à ajouter la ligne \texttt{Logged(Logged\&\&) = default;}.

Cette seule ligne ne suffit pas tout à fait. Déclarer un constructeur de déplacement supprime les opérations de copie implicites~; il faut donc ajouter \texttt{Logged(const Logged\&) = default;} ainsi que les deux affectations, si on les souhaite, et \texttt{Logged() = default;} également.

Rien ne le signale, puisqu'il s'agit d'un changement silencieux de performance~: \texttt{std::move} compile toujours, et une copie est une façon valide de satisfaire une demande de déplacement. Ainsi, un membre comme une longue chaîne ou un grand vecteur est désormais dupliqué à chaque déplacement.

La leçon porte sur ce qui semble anodin. Un destructeur ajouté pour journaliser quelque chose semble gratuit, et il peut modifier la manière dont tous les objets de la classe se déplacent. Si l'on veut voir les déplacements, il faut les compter, comme le fait le listing.

\iqwrong{«~Un destructeur vide revient à n'en avoir aucun.~» Son corps ne fait rien, mais il est déclaré par l'utilisateur, et c'est cela que le langage examine~: les déplacements disparaissent, et un type trivialement copiable cesse de l'être, ce qui change aussi la manière dont il est passé en argument.}
\iqfollow{Que fait \texttt{= default} sur le destructeur~?}{Il reste déclaré par l'utilisateur, donc il supprime les déplacements exactement de la même manière.}
\end{solution}

\begin{solution}{iq:iv3:cpp:15}
Il affiche \texttt{0 copies, 7 moves} puis \texttt{7 copies, 0 moves}. La capacité vaut successivement 1, 2, 4, puis 8, si bien que les cinq insertions déplacent 1, 2 et 4 anciens éléments. Un vecteur ne les déplace que si le déplacement ne peut pas lever d'exception. Sinon il copie, afin de conserver la garantie forte d'exception de \texttt{push\_back}.

La garantie dit que si l'insertion lève une exception, le vecteur reste tel qu'il était. Avec des copies, c'est facile à tenir~: si la troisième copie lève une exception, on détruit les nouvelles copies, on libère le nouveau bloc, et les anciens éléments restent intacts. Avec un déplacement pouvant lever une exception, les premiers éléments de l'ancien bloc sont déjà dans un état déplacé quand le troisième déplacement échoue, et les déplacer en arrière pourrait lever une exception aussi, si bien qu'il n'existe aucun retour possible.

Le vecteur utilise donc \texttt{std::move\_if\_noexcept}, qui ne donne un déplacement que si celui-ci est \texttt{noexcept} ou si le type ne peut pas être copié du tout. Un type non copiable dont le déplacement peut lever une exception se déplace tout de même, et la garantie est alors abandonnée.

Un constructeur de déplacement qui ne lève jamais d'exception doit donc être déclaré \texttt{noexcept}. Sans cela, les déplacements écrits peuvent ne jamais s'exécuter dans un vecteur. Un déplacement généré par le compilateur est \texttt{noexcept} de lui-même lorsque les déplacements de tous les membres le sont.

\iqwrong{«~Le compilateur voit que ce déplacement ne peut pas lever d'exception.~» Il vérifie la déclaration. Un corps en ligne qui ne fait rien est tout de même considéré comme pouvant lever une exception ici.}
\iqfollow{Comment vérifiez-vous un type~?}{Avec un \texttt{static\_assert} placé à côté du type, sur le trait nothrow-move de \texttt{<type\_traits>}.}
\end{solution}

\begin{solution}{iq:iv3:cpp:16}
Il affiche \texttt{0 40}, \texttt{3}, \texttt{1 1}, \texttt{1 5} et \texttt{new!}. La norme garantit le 40, le 3, le \texttt{1 1}, le \texttt{1 5} et la dernière ligne. La taille 0 d'une chaîne déplacée est ce que fait ici libstdc++, et la norme la qualifie seulement de «~valide mais non spécifiée~».

\texttt{std::move} ne fait rien à l'exécution. C'est un transtypage vers une référence rvalue, qui permet de choisir un constructeur de déplacement. Les deux dernières lignes du listing le montrent~: \texttt{r} n'est qu'un autre nom pour \texttt{a}, et \texttt{a} se modifie par son intermédiaire, sans que rien ne soit déplacé.

La suite découle de ce qui est déplacé. Un \texttt{int} n'a aucune ressource, si bien que son déplacement est une copie, et \texttt{i} reste à 1. Un \texttt{unique\_ptr} est garanti nul après avoir été déplacé, puisque sa spécification le dit explicitement.

Sur une chaîne déplacée, on peut faire tout ce qui n'a pas de précondition~: lui affecter une valeur, comme le montre la deuxième ligne, la détruire, ou demander sa taille. On ne peut pas supposer une valeur, si bien que \texttt{a.front()} n'est pas sûr sans vérification.

\iqwrong{«~Après \texttt{std::move(a)}, \texttt{a} est vide.~» Le transtypage seul ne change rien, et même un véritable déplacement ne garantit pas une chaîne vide.}
\iqfollow{Que se passe-t-il avec une chaîne courte~?}{Ses caractères se trouvent dans l'objet lui-même, si bien que le déplacement les copie. Rien n'y est moins coûteux.}
\end{solution}

\begin{solution}{iq:iv3:cpp:17}
Il affiche \texttt{plain: 0 copies, 0 moves}, \texttt{forced: 0 copies, 1 moves} et \texttt{either: 1 copies, 0 moves}. \texttt{std::move} appliqué au retour empêche l'élision, car l'expression renvoyée n'est plus le simple nom local. Dans \texttt{either}, l'expression renvoyée \texttt{c ? p : q} n'est pas le nom d'une variable locale, si bien que le résultat est copié.

Renvoyer une variable locale par son nom permet au compilateur de la construire directement à l'emplacement de l'appelant, ce que fait g++ ici. C'est permis mais non obligatoire. S'il s'y refuse, le langage tente tout de même un déplacement d'abord, si bien que \texttt{return t;} n'a jamais besoin de \texttt{std::move}. Le transtypage n'apporte donc rien et coûte un déplacement. \texttt{-Wpessimizing-move}, présent dans \texttt{-Wall}, le signale, et le listing le désactive sur une ligne non montrée.

L'expression conditionnelle est un cas différent. Ce n'est pas un nom, donc aucun déplacement implicite ne s'applique. \texttt{if (c) return p; return q;} déplacerait, ou pourrait élider.

\texttt{make} fonctionne pour une raison différente. Depuis C++17, une fonction qui renvoie directement une prvalue n'a aucun temporaire à copier~: \texttt{Pinned\{\}} initialise \texttt{p} lui-même. Il s'agit donc d'une règle du langage, et aucun constructeur de copie n'a besoin d'exister.

\iqwrong{«~Écrire toujours \texttt{std::move} au retour, pour être sûr que cela déplace.~» C'est le seul cas où il nuit. Il ne sert alors que pour renvoyer un membre, ou un paramètre qui est une référence rvalue, dans les normes antérieures.}
\end{solution}

\begin{solution}{iq:iv3:cpp:18}
Il affiche \texttt{copy}, \texttt{copy}, \texttt{copy}, \texttt{move}, \texttt{move}. Dans \texttt{bad}, \texttt{x} est une variable nommée, et toute variable nommée est une lvalue, si bien que \texttt{take(x)} choisit toujours la surcharge \texttt{const std::string\&}.

\texttt{std::forward<T>(x)} rétablit ce que l'appelant a transmis~: une lvalue reste une lvalue, et une rvalue redevient une rvalue. Il peut le faire parce que \texttt{T} porte cette information~: pour une lvalue de l'appelant, \texttt{T} vaut \texttt{std::string\&}, et pour une rvalue, c'est \texttt{std::string}.

\texttt{std::move} fait de tout une rvalue, sans condition. Avec lui dans \texttt{good}, l'appel \texttt{good(s)} déplacerait le contenu de la variable \texttt{s} de l'appelant, la vidant alors que l'appelant ne l'avait jamais abandonnée. C'est un bogue désagréable, puisque le code de l'appelant ne montre aucun \texttt{std::move}.

La règle est donc simple~: \texttt{std::forward} s'emploie avec une référence de transmission, et \texttt{std::move} avec une référence rvalue. Ne transmettre un argument qu'une fois au plus, puisqu'un second usage lirait un objet dans un état déplacé.

\iqwrong{«~\texttt{x} a le type \texttt{T\&\&}, donc c'est une rvalue.~» Le type de la variable et la catégorie de l'expression sont deux choses différentes. Un nom est toujours une lvalue.}
\iqfollow{Qu'est-ce qui pose problème avec un pointeur de fonction ou une liste entre accolades~?}{La transmission parfaite n'est pas parfaite. \texttt{good(\{1, 2\})} échoue, car une liste entre accolades n'a pas de type à partir duquel déduire \texttt{T}. Un pointeur de fonction se transmet correctement, mais le nom d'une fonction surchargée ne le fait pas~: lui non plus n'a de type unique à partir duquel déduire \texttt{T}.}
\end{solution}

\begin{solution}{iq:iv3:cpp:19}
Il affiche \texttt{T is a reference}, \texttt{T is a value} puis \texttt{T is a reference}. Seuls \texttt{T\&\&} dans \texttt{f} et \texttt{auto\&\&} sont des références de transmission, puisque chacun a exactement la forme d'un nom déduit suivi de \texttt{\&\&}. \texttt{g(l)} ne compile pas, car \texttt{std::vector<T>\&\&} est une simple référence rvalue et \texttt{l} est une lvalue.

Avec \texttt{f(i)}, \texttt{i} est une lvalue, si bien qu'une règle de déduction particulière donne à \texttt{T} le type \texttt{int\&}. Le type du paramètre devient \texttt{int\& \&\&}. On ne peut pas l'écrire directement, mais il peut apparaître à partir d'un modèle ou d'un alias, et la réduction des références le simplifie~: si l'une des deux parties est \texttt{\&}, le résultat est \texttt{\&}, et seul \texttt{\&\& \&\&} donne \texttt{\&\&}. Le paramètre vaut donc \texttt{int\&}. Avec \texttt{f(1)}, \texttt{T} vaut \texttt{int}, et le paramètre est \texttt{int\&\&}. Avec un \texttt{const int}, \texttt{T} vaut \texttt{const int\&}.

Toute adjonction d'un qualificatif fait perdre à la référence son statut de référence de transmission, comme dans \texttt{const T\&\&} ou dans \texttt{std::vector<T>\&\&}. Un \texttt{T\&\&} figurant dans un membre d'une classe patron, où \texttt{T} a été fixé lors de l'instanciation de la classe, n'en est pas une non plus, car aucune déduction n'a lieu à cet appel.

Ainsi, \texttt{auto\&\& a = i} fait de \texttt{a} un \texttt{int\&}, et \texttt{auto\&\& b = 1} en fait un \texttt{int\&\&}, ce que le listing vérifie avec \texttt{static\_assert}.

\iqwrong{«~\texttt{\&\&} désigne toujours une référence rvalue.~» Seulement là où rien n'est déduit.}
\end{solution}

\begin{solution}{iq:iv3:cpp:20}
\texttt{make()} renvoie un temporaire, détruit à la fin de la déclaration, si bien que \texttt{v} pointe vers une zone mémoire libérée. AddressSanitizer signale \texttt{heap-use-after-free} à la ligne qui lit \texttt{v}. Une chaîne de 5~caractères serait conservée à l'intérieur de l'objet \texttt{std::string}, par l'optimisation des chaînes courtes, de sorte qu'aucun tas n'est impliqué~; la vue pendrait alors dans la pile devenue invalide.

Un \texttt{string\_view} est un pointeur et une longueur, sans en être propriétaire. Lier une référence \texttt{const std::string\&} à un temporaire prolonge la vie de celui-ci. Une vue n'a pas cet effet, puisque sa conversion passe par un opérateur, et le temporaire meurt comme d'habitude. Il faut donc garder la chaîne dans une variable nommée, et prendre la vue à partir de celle-ci.

En paramètre, une vue est excellente~: \texttt{void parse(std::string\_view)} accepte des littéraux et des chaînes sans copie, et l'argument de l'appelant survit à l'appel. Comme membre ou comme valeur de retour, elle est mauvaise, car l'objet peut survivre à ce qu'il regarde, et rien dans le type ne le signale. Si une classe doit conserver le texte, elle doit contenir un \texttt{std::string}.

\iqwrong{«~Il affiche le bon résultat, donc c'est bon.~» Une zone mémoire libérée conserve souvent les anciens octets. Le programme a un comportement indéfini, et il passe jusqu'à ce qu'une allocation réutilise le bloc.}
\iqfollow{Est-ce encore vrai pour un \texttt{std::span}~?}{Oui. C'est un pointeur et une taille, et il obéit à la même règle de durée de vie.}
\end{solution}

\begin{solution}{iq:iv3:cpp:21}
Il affiche \texttt{4 12}, puis \texttt{4}. \texttt{strlen} lit jusqu'à un octet nul. \texttt{sym.data()} pointe dans le littéral entier, et le seul octet nul se trouve à sa fin, si bien que la fonction C voit \texttt{ESZ6,4501.25}, soit 12~caractères, alors que la vue n'en a que 4.

Le bogue tient donc à une hypothèse~: une vue est un pointeur plus une longueur, et aucun octet nul n'est tenu de suivre la vue. Une fonction C n'a aucun moyen de recevoir la longueur, elle continue donc de lire. La sous-chaîne d'un texte plus long se termine au milieu d'autres caractères, et c'est le cas usuel pour un analyseur syntaxique.

Ici, ce n'est pas un comportement indéfini, puisque le littéral possède bien un octet nul à la fin. C'est un résultat faux. Avec une vue sur un tampon non terminé par un octet nul, comme des octets issus d'un socket, \texttt{strlen} dépasserait la fin, ce qui constitue un comportement indéfini.

Le programme passe les tests, car une vue sur un littéral entier se termine bien par un octet nul, si bien que le cas simple fonctionne. Seule une véritable sous-chaîne le révèle. Le correctif pour un appel C consiste à faire une copie terminée, \texttt{std::string(sym).c\_str()}, comme à la deuxième ligne, et à garder cette chaîne vivante pendant l'appel. Ou utiliser une fonction qui prend une longueur.

\iqwrong{«~\texttt{string\_view} possède bien \texttt{c\_str()}, non~?~» Il n'a pas ce membre, justement parce qu'il ne peut pas garantir la terminaison.}
\end{solution}

\begin{solution}{iq:iv3:cpp:22}
Il affiche \texttt{2 2}, et \texttt{freed} ne s'affiche jamais. Chaque nœud est retenu par une variable locale et par l'autre nœud. Quand les variables locales disparaissent, chaque compteur retombe à 1, si bien qu'aucun n'atteint 0. LeakSanitizer signale 64~octets dans 2~allocations. Le \Cref{lst:iv3:cp:cyclefixed} brise le cycle avec un seul lien faible, et le test montre que les deux destructeurs s'exécutent.
\omcode{code/interviews-3/03-cpp/cpp/cycle_fixed.cpp}{1}{22}{Un lien fort et un lien faible.}\label{lst:iv3:cp:cyclefixed}

Le comptage de références ne voit pas un cycle, car il ne connaît que des compteurs locaux, et plus rien à l'extérieur ne désigne la paire. Un ramasse-miettes qui parcourt à partir des racines les trouverait, mais C++ n'en a pas.

Le choix porte sur le lien qui désigne la propriété. Ici \texttt{next} possède et \texttt{prev} observe. Un \texttt{weak\_ptr} ne maintient pas l'objet en vie et ne peut pas être déréférencé directement, il faut donc appeler \texttt{lock()}. On obtient alors un \texttt{shared\_ptr}, ou un pointeur nul si l'objet a disparu, et c'est pourquoi le listing le teste. Tant qu'on tient le résultat, l'objet ne peut pas mourir sans prévenir. C'est le compteur de 2 de la ligne \texttt{alive 2}.

Après le bloc, \texttt{expired()} vaut vrai, et \texttt{lock()} renvoie un pointeur nul.

\iqwrong{«~\texttt{use\_count} me dit s'il est sûr de l'utiliser.~» Avec plusieurs threads, sa valeur peut être périmée dès qu'elle est renvoyée. \texttt{lock()} est la vérification atomique.}
\iqfollow{Qu'apporte un pointeur brut par rapport à un \texttt{weak\_ptr} ici~?}{Il ne coûte rien, et convient si le propriétaire survivra sûrement à ce pointeur. Un \texttt{weak\_ptr} sert lorsque ce n'est pas certain.}
\end{solution}

\begin{solution}{iq:iv3:cpp:23}
Le \Cref{lst:iv3:cp:unique} est une telle classe, aussi grande qu'un pointeur brut, soit 8~octets sur x86-64. L'auto-déplacement est sûr, car \texttt{release} vide d'abord l'objet et \texttt{reset} reprend le même pointeur. Elle se comporte comme \texttt{std::unique\_ptr} sur 20\,000 opérations aléatoires, sans fuite sous AddressSanitizer.
\omcode{code/interviews-3/03-cpp/cpp/unique.hpp}{4}{28}{Un petit pointeur possédant sa ressource.}\label{lst:iv3:cp:unique}

La propriété signifie un seul propriétaire, donc la copie est supprimée, et un déplacement transmet le pointeur et laisse un pointeur nul. Les deux sont \texttt{noexcept}, si bien qu'un vecteur de ces objets se déplace lors de sa croissance.

\texttt{reset} enregistre d'abord le nouveau pointeur, car le destructeur de l'ancien objet peut revenir ici. Si un nœud possède un enfant au moyen d'un \texttt{Unique}, et que le destructeur de cet enfant finit par appeler \texttt{reset} sur le même pointeur, une version qui libère d'abord l'ancien supprimerait l'objet deux fois. \texttt{std::exchange} laisse le membre cohérent avant que tout code extérieur ne s'exécute.

Le mot-clé \texttt{explicit} sur le constructeur compte. Sans lui, un pointeur brut se convertirait silencieusement, et le passer à une fonction qui prend un \texttt{Unique} le détruirait à la fin de l'appel.

\texttt{delete nullptr} est défini et ne fait rien, si bien qu'aucun test n'est nécessaire dans le destructeur. Le type doit être complet à ce point, sans quoi le compilateur ne peut pas appeler son destructeur, ce qui compte lorsque la classe est utilisée avec une déclaration anticipée, comme dans l'idiome pimpl.

\iqwrong{\texttt{delete p\_; p\_ = o.p\_; o.p\_ = nullptr;} dans l'affectation par déplacement, sans aucun contrôle. À l'auto-déplacement, l'objet est supprimé puis le pointeur se retrouve nul~: la valeur est détruite sans bruit.}
\iqfollow{Qu'ajoute la version réelle~?}{Un type de suppresseur, la prise en charge des tableaux, et la conversion depuis un pointeur vers une classe dérivée.}
\end{solution}

\begin{solution}{iq:iv3:cpp:24}
Les deux affichent \texttt{destroyed} quand le dernier \texttt{shared\_ptr} disparaît. Avec \texttt{new}, il y a alors 1 libération, puis 2 après la disparition du \texttt{weak\_ptr}. Avec \texttt{make\_shared}, il n'y en a alors aucune, et 1 après. Dans le second cas, le mégaoctet reste alloué jusqu'à ce que le dernier \texttt{weak\_ptr} ait disparu.

\texttt{make\_shared} place l'objet et le bloc de contrôle dans une seule allocation. Le bloc de contrôle doit survivre à l'objet, car un \texttt{weak\_ptr} a besoin de ses compteurs pour répondre à \texttt{expired()}. Le destructeur s'exécute quand le compteur fort atteint 0, mais la mémoire ne forme qu'un seul bloc, qui ne peut donc pas être rendu en partie. Avec \texttt{new}, il y a deux blocs, et celui de l'objet est libéré en premier.

Il faut donc éviter \texttt{make\_shared} pour un grand objet observé par des pointeurs faibles de longue durée, comme un cache qui conserve des références faibles. Un destructeur qui libère la mémoire d'un membre aide d'ailleurs, puisque seuls les octets propres à l'objet subsistent.

Sinon, \texttt{make\_shared} est le choix habituel. Une seule allocation au lieu de deux signifie moins de temps et moins de surcharge, et l'objet et ses compteurs sont proches l'un de l'autre, ce qui favorise le cache. On n'écrit pas non plus \texttt{new}. Il ne peut pas recevoir de suppresseur personnalisé, et il exige un constructeur public.

\iqwrong{«~Le destructeur a été exécuté, donc la mémoire est libérée.~» La destruction et la libération sont deux étapes, qui ici ont lieu à des moments différents.}
\end{solution}

\begin{solution}{iq:iv3:cpp:25}
\texttt{a} et \texttt{b} créent chacun leur propre bloc de contrôle pour un même objet, et chaque compteur vaut 1, si bien qu'ils ne se connaissent pas. Le dernier à disparaître supprime de nouveau l'objet, et AddressSanitizer signale \texttt{attempting double-free}.

La propriété n'est partagée qu'en copiant un \texttt{shared\_ptr}. Construire à partir d'un pointeur brut démarre toujours un nouveau compteur. Il faut donc tenir le pointeur brut à l'écart~: \texttt{auto a = std::make\_shared<Order>(); auto b = a;} donne des compteurs qui valent alors 2.

\texttt{std::shared\_ptr<Order>(this)} démarre un second compteur sur un objet déjà possédé, ce qui est le même bogue, avec les deux propriétaires encore plus éloignés l'un de l'autre.

La bonne façon, dans le \cref{lst:iv3:cp:sharethis}, consiste à dériver de \texttt{std::enable\_shared\_from\_this} et à appeler \texttt{shared\_from\_this()}. Quand un \texttt{shared\_ptr} prend la propriété, il conserve dans cette base un \texttt{weak\_ptr} vers lui-même, et l'appel le verrouille, ce qui partage le même bloc de contrôle.
\enlargethispage{-3\baselineskip}%
\omcode{code/interviews-3/03-cpp/cpp/share_this.cpp}{1}{18}{Donner un pointeur vers soi-même.}\label{lst:iv3:cp:sharethis}

Il faut donc que l'objet soit déjà possédé par un \texttt{shared\_ptr}. Si ce n'est pas le cas, comme avec \texttt{plain}, le pointeur faible est vide, et depuis C++17 l'appel lève \texttt{std::bad\_weak\_ptr}. Il n'est pas non plus possible dans le constructeur, puisqu'aucun \texttt{shared\_ptr} n'a encore pris la propriété.

\iqwrong{«~Les compteurs sont atomiques, donc deux propriétaires sont sûrs.~» L'atomicité porte sur un compteur vu par plusieurs threads. Il y a ici deux compteurs.}
\end{solution}

\begin{solution}{iq:iv3:cpp:26}
\texttt{push\_back} a besoin d'une capacité supérieure à celle de 3 du vecteur, si bien qu'il déplace les éléments vers un nouveau bloc et libère l'ancien. \texttt{first} pointe encore dans le bloc libéré. AddressSanitizer signale \texttt{heap-use-after-free} à la ligne \texttt{first = 9}.

Toute opération qui réalloue invalide toutes les références, tous les pointeurs et tous les itérateurs pointant dans un vecteur. Si la capacité n'est pas dépassée, ceux qui précèdent le point d'insertion restent valides. Ce bogue dépend donc de la capacité au moment considéré, et c'est pour cela qu'il se cache~: un vecteur qui avait déjà de la place libre s'exécute correctement.

Deux correctifs. Appeler d'abord \texttt{v.reserve(4)}, et la référence reste alors valide, comme le confirme une exécution propre sous le sanitizer. C'est fragile, car la réservation doit couvrir chaque insertion ultérieure. Ou conserver un indice, \texttt{std::size\_t first = 0}, et écrire \texttt{v[first] = 9}, ce qui survit à toute réallocation, et c'est la réponse habituelle.

\texttt{std::deque} suit d'autres règles. \texttt{push\_back} et \texttt{push\_front} invalident les itérateurs mais non les références aux éléments, si bien que ce code exact est correct sur un deque. Une insertion au milieu invalide les deux.

\iqwrong{«~Il a affiché 1 ici, donc l'écriture a été perdue, et c'est le seul dommage.~» Écrire dans une mémoire libérée peut corrompre un autre objet ou les données propres à l'allocateur, loin de cette ligne.}
\iqfollow{Et une range-for qui ajoute des éléments au vecteur qu'elle parcourt~?}{Même bogue~: ses itérateurs sont invalidés dès la première réallocation.}
\end{solution}

\begin{solution}{iq:iv3:cpp:27}
Il affiche \texttt{0 1 1}. \texttt{auto y = w[0]} copie un \texttt{int}. \texttt{vector<bool>} stocke des bits, et un bit n'a pas d'adresse, si bien que \texttt{v[0]} renvoie un petit objet proxy, et \texttt{auto x} est ce proxy. Affecter \texttt{x} écrit directement dans le bit, si bien que \texttt{v[0]} vaut vrai.

\texttt{bool\& b = v[0];} ne compile pas, puisque le proxy n'est pas un \texttt{bool}. Ni \texttt{auto\& b = v[0];}, car le proxy est un temporaire. \texttt{bool b = v[0];} compile, et il fait une copie indépendante, ce que \texttt{auto} ne faisait pas.

Deux indices différents de \texttt{vector<bool>} peuvent partager un octet, si bien qu'écrire dans l'un revient à lire et écrire cet octet, et la norme ne garantit pas la sécurité~: \texttt{vector<bool>} est l'exception nommée à la règle selon laquelle des éléments différents d'un conteneur peuvent être écrits simultanément. ThreadSanitizer signale une course aux données (data race) dans le test de \texttt{vector<bool>}, et aucune pour \texttt{vector<char>}.

Donc \texttt{vector<bool>} n'est pas vraiment un conteneur de \texttt{bool}~: un \texttt{bool*} ne peut pas pointer vers un élément, et le code générique qui prend un \texttt{T\&} échoue. Utiliser \texttt{vector<char>} lorsqu'on veut le comportement habituel, et \texttt{std::bitset} ou \texttt{vector<bool>} lorsqu'on veut le gain de mémoire d'un bit par élément.

\iqwrong{«~Ça économise de la mémoire, donc c'est toujours mieux.~» Chaque accès demande un décalage et un masque, et la règle sur les courses est différente. Dans une boucle serrée, les octets peuvent être plus rapides.}
\end{solution}

\begin{solution}{iq:iv3:cpp:28}
Il affiche \texttt{1 1}, \texttt{1} et \texttt{0}. La norme garantit le second nombre de la première ligne et la seconde ligne~: les références et pointeurs vers les éléments d'un \texttt{unordered\_map} et d'un \texttt{map} restent valides lors des insertions. Le premier nombre indique que la table a été rehachée, ce qui dépend du nombre initial de seaux choisi par la bibliothèque. Le \texttt{0} indique que le vecteur a été déplacé, et aucune garantie ne couvre cette valeur non plus.

Un rehachage invalide les itérateurs d'un \texttt{unordered\_map}, et non les références ou pointeurs vers ses éléments. Les éléments vivent dans des nœuds alloués séparément dans libstdc++, et la règle de la norme impose une telle conception. C'est pourquoi \texttt{unordered\_map} ne peut pas être une table plate à adressage ouvert, et c'est le prix de cette garantie.

Pour les pointeurs vers des ordres conservés dans un conteneur~: \texttt{std::map}, \texttt{std::unordered\_map}, \texttt{std::list} et \texttt{std::deque} (aux extrémités) le permettent tous. \texttt{std::vector} ne le permet pas, sauf si la capacité est fixée à l'avance. Un carnet d'ordres conserve donc souvent les ordres dans un pool aux adresses stables, avec un index par-dessus, ou stocke des indices entiers, ce qui exploite aussi mieux le cache.

L'effacement est une question à part. Il n'invalide que les références vers l'élément effacé, sur \texttt{map} et \texttt{unordered\_map}.

\iqwrong{«~Un \texttt{unordered\_map} déplace tout lors du rehachage, donc les pointeurs meurent.~» Cela le confond avec un vecteur. Ses nœuds restent en place, et seul le tableau des seaux est reconstruit.}
\iqfollow{Que sacrifie une table de hachage plate pour aller plus vite~?}{Exactement ceci~: ses éléments se déplacent lors de la croissance, si bien que les pointeurs ne survivent pas.}
\end{solution}

\begin{solution}{iq:iv3:cpp:29}
Il affiche \texttt{1}, \texttt{1} et \texttt{1000000}. Le premier compte vient de \texttt{push\_back(T(1))}, qui construit un temporaire et le déplace à l'intérieur. \texttt{emplace\_back(2)} construit sur place, sans déplacement, si bien que le compte reste à 1. Le \texttt{push\_back} d'un simple 1000000 ne compile pas, et c'est tout l'intérêt de l'exemple.

\texttt{emplace\_back} est utile lorsqu'on passe des arguments de constructeur et qu'un déplacement coûte cher, ou n'est pas disponible du tout. Si l'on dispose déjà d'un objet, \texttt{push\_back(std::move(x))} et \texttt{emplace\_back(std::move(x))} font le même travail. Pour un type léger, comme une petite structure d'entiers, la différence est quasi nulle une fois optimisé.

Le danger est qu'\texttt{emplace\_back} puisse appeler des constructeurs \texttt{explicit}. Le constructeur de vecteur qui prend une taille est \texttt{explicit}, afin qu'un nombre ne devienne pas silencieusement un vecteur. Ainsi, l'\texttt{emplace\_back} du listing crée une ligne d'un million de zéros, alors que l'auteur voulait peut-être une ligne contenant un seul nombre. \texttt{push\_back} exige une conversion implicite, et il n'y en a pas, si bien que l'erreur est détectée. Avec des accolades, \texttt{push\_back(\{1000000\})} donne la ligne d'un seul élément.

Il faut donc utiliser \texttt{push\_back} par défaut, pour sa vérification, et \texttt{emplace\_back} lorsqu'on construit vraiment sur place.

\iqwrong{«~\texttt{emplace\_back} est toujours plus rapide, donc il faut l'utiliser partout.~» Il peut être équivalent, il peut être plus lent à compiler, et il désactive une vérification de sûreté.}
\end{solution}

\begin{solution}{iq:iv3:cpp:30}
AddressSanitizer signale un \texttt{heap-buffer-overflow}, une lecture au-delà de la fin, dans \texttt{\_\_unguarded\_partition}, à l'intérieur du tri. Le comparateur a violé l'exigence d'un ordre faible strict~: \texttt{comp(a, a)} doit valoir faux, et avec \texttt{<=} il vaut vrai. Violer cette exigence rend l'appel indéfini, si bien que ce plantage n'est qu'un résultat parmi d'autres autorisés.

L'implémentation s'appuie sur cette garantie. L'étape de partition de libstdc++ parcourt vers l'avant \emph{sans contrôle de bornes}, en comptant sur le fait que la médiane de trois a laissé dans l'intervalle un élément qui n'est pas inférieur au pivot. Avec \texttt{<=}, un élément égal compte aussi comme inférieur, si bien que le parcours ne s'arrête pas là et dépasse la fin du tableau.

Le test passe sur les petites entrées, car libstdc++ utilise un tri par insertion pour les intervalles courts, sans cette étape~: le seuil est de 16 éléments dans cette version. Il passe aussi sur des valeurs distinctes, car \texttt{<=} et \texttt{<} donnent alors les mêmes réponses. Seules les nombreuses clés égales le révèlent, ce qui arrive la première fois qu'un vrai jeu de données arrive avec des cours répétés.

Pour deux champs, il faut comparer un tuple~: \texttt{std::tie(a.px, a.t) < std::tie(b.px, b.t)}, qui donne la bonne logique, ou, en C++20, un \texttt{operator<=>} défini par défaut. L'écrire à la main invite à l'erreur~: \texttt{a.px < b.px || a.t < b.t} n'est pas non plus un ordre faible strict.

\iqwrong{«~Il trie un peu bizarrement, mais le pire cas est un ordre faux.~» Un comparateur défaillant peut lire et écrire en dehors du tableau.}
\iqfollow{Quelles autres structures ont besoin de cela~?}{\texttt{std::map}, \texttt{std::set} et les tas reposent sur la même règle.}
\end{solution}

\begin{solution}{iq:iv3:cpp:31}
Il affiche \texttt{0 1}~: le premier test ne détecte pas le dépassement, et le compilateur le détecte. Le dépassement d'entier signé est indéfini, si bien que le compilateur peut supposer que \texttt{a + b} ne déborde pas. Dans ce cas, \texttt{a + b < a} équivaut à \texttt{b < 0}, et c'est ce qui est compilé.

Avec $b = 1$, ce test est faux, quelle que soit la valeur de $a$, si bien que le contrôle ne peut jamais se déclencher. Le code se lit comme si le dépassement se repliait, ce qui est le cas sur le matériel, mais l'optimiseur raisonne à partir des règles du langage, et non de la machine.

Sous UBSan, l'exécution signale \texttt{signed integer overflow: 2147483647 + 1} et affiche \texttt{1 1}, car la version instrumentée effectue l'addition et elle se replie. Avec \texttt{-fwrapv}, le dépassement est défini comme un repli, si bien que le contrôle fonctionne, et le programme affiche aussi \texttt{1 1} ici. Ces deux options changent ce qui est demandé, et aucune ne rend ce code correct en C++.

Le test portable vérifie avant d'additionner~: pour $b > 0$, \texttt{a > INT\_MAX - b}, et pour $b < 0$, \texttt{a < INT\_MIN - b}. Aucun des deux membres ne peut déborder. \texttt{\_\_builtin\_add\_overflow} est une extension du compilateur dans g++ et Clang qui fait la même chose, et ici elle lit l'indicateur de dépassement du processeur avec un \texttt{seto}, comme le confirme notre test.

\iqwrong{«~Ça doit être un bogue du compilateur, puisque le matériel se replie.~» C'est le programme qui a le bogue. g++ 11.4 applique cette réécriture même en \texttt{-O0}~: le listing affiche aussi \texttt{0 1} là.}
\end{solution}

\begin{solution}{iq:iv3:cpp:32}
Il affiche \texttt{-2147483648}, \texttt{2}, \texttt{-2147483648} et \texttt{4294967295}. Le premier et le dernier sont définis par la norme, et le deuxième et le troisième sont indéfinis, si bien qu'UBSan signale ces deux-là.

\texttt{1 <{}< 31} est défini depuis C++20, car les entiers signés sont en complément à deux et les bits vont là où ils vont, ce qui donne ici \texttt{INT\_MIN}. Avant C++20, sa valeur était définie par l'implémentation. \texttt{u - 1} sur un 0 non signé se replie vers $2^{32} - 1$, ce qui est toujours défini.

\texttt{1 <{}< 33} sur un type de 32 bits a un nombre de décalages égal ou supérieur à la largeur, ce qui reste indéfini en C++20. Il affiche 2 parce que l'instruction de décalage de x86 n'utilise que les 5 bits de poids faible du nombre de décalages, si bien que 33 agit comme 1. C'est le comportement de cette machine, et rien n'oblige une autre machine ou un autre compilateur à s'y conformer.

\texttt{std::abs(INT\_MIN)} n'a pas de réponse dans \texttt{int}, car $2^{31}$ ne tient pas dans ce type, et le résultat est donc indéfini. Sur cette machine, la négation se replie, ce qui laisse une valeur absolue négative. Tout code qui suppose qu'une valeur absolue est positive ou nulle échoue alors, par exemple en l'utilisant comme indice modulo la taille d'une table.

Les rapports d'UBSan indiquent \texttt{shift exponent 33 is too large} et \texttt{negation of -2147483648}.

\iqwrong{«~Indéfini veut dire plantage.~» Ici, les deux lignes indéfinies affichent des nombres plausibles. C'est la manière habituelle, et c'est pourquoi il faut le sanitizer.}
\iqfollow{Et \texttt{INT\_MIN / -1}~?}{Également indéfini, car le vrai résultat ne tient pas dans le type, et sur x86-64 l'instruction de division déclenche une trappe. Notre test voit \texttt{SIGFPE}.}
\end{solution}

\begin{solution}{iq:iv3:cpp:33}
\texttt{one} et \texttt{two} sont sûrs, et \texttt{three} pend dans le vide. Lier directement un temporaire à une référence \texttt{const} prolonge la durée de vie du temporaire jusqu'à celle de la référence, et cela vaut aussi pour un membre du temporaire, qui est \texttt{two}. AddressSanitizer signale un \texttt{stack-use-after-scope} à la dernière lecture.

Cette prolongation est une propriété de l'initialisation, et le compilateur l'applique là où il voit le temporaire être lié. Il ne regarde pas à l'intérieur d'une fonction. Pour \texttt{same(make())}, le temporaire se lie au paramètre \texttt{s}. Cela dure jusqu'à la fin de l'expression complète, et la fonction renvoie une référence vers lui. Le compilateur ne peut pas savoir, en général, que la référence renvoyée désigne cet argument, si bien que \texttt{three} est lié à un objet sur le point de disparaître.

Il en va de même pour toute fonction qui renvoie une référence vers son argument, ce qui inclut \texttt{std::max}, \texttt{std::min} et \texttt{std::move}. Une fonction membre qui renvoie une référence vers un membre d'un temporaire se comporte de la même façon~: \texttt{Holder\{make()\}.get\_name()} pendrait dans le vide, alors que l'accès direct \texttt{.name} ne le faisait pas.

Donc, dans le doute, il faut prendre une valeur. \texttt{std::string three = same(make());} copie pendant que le temporaire est encore vivant, et ne coûte qu'une copie.

\iqwrong{«~Une référence \texttt{const} garde toujours un temporaire en vie.~» Seulement celle qui y est liée directement, dans la déclaration.}
\iqfollow{La forme range-for est-elle sûre~?}{\texttt{for (auto\& x : make\_vector())} l'est, car l'intervalle est lié directement. Une chaîne comme \texttt{make\_holder().items()} ne l'est pas, avant C++23.}
\end{solution}

\begin{solution}{iq:iv3:cpp:34}
Il affiche \texttt{Derived 1}, \texttt{Derived 2}, \texttt{Base::g(int)} et \texttt{Derived::g(double)}. Les arguments par défaut sont choisis à la compilation à partir du type statique, et le corps de la fonction à l'exécution à partir de l'objet. Les noms déclarés dans une classe dérivée masquent toutes les surcharges de la classe de base portant ce nom.

Ainsi, \texttt{b.f()} prend son argument par défaut dans \texttt{Base}, car \texttt{b} est déclaré \texttt{Base\&}, puis il se répartit vers \texttt{Derived::f}, qui reçoit 1. Deux valeurs par défaut pour une même fonction constituent un piège sans avertissement. La solution n'est pas de redéfinir les valeurs par défaut dans une redéfinition. Mieux vaut placer la valeur par défaut dans une fonction non virtuelle qui appelle une fonction virtuelle sans valeur par défaut.

La recherche pour \texttt{d.g(1)} s'arrête au premier espace de noms qui contient le nom, c'est-à-dire \texttt{Derived}. \texttt{Base::g(int)} n'est même pas considérée, si bien que 1 se convertit en \texttt{double}. La correction consiste à écrire \texttt{using Base::g;} dans \texttt{Derived}, ce qui fait entrer les surcharges de la base dans l'ensemble, et alors \texttt{d.g(1)} choisit la correspondance exacte.

Une option du compilateur peut avertir d'un cas précis~: \texttt{-Woverloaded-virtual} signale une fonction qui masque une fonction virtuelle de la base. Ici, \texttt{g} n'est pas virtuelle, si bien que cette option ne sert à rien, et le listing compile sans avertissement. Il faut donc le chercher à la relecture.

\iqwrong{«~\texttt{override} aurait détecté cela.~» Il vérifie que \texttt{f} redéfinit quelque chose. Il ne dit rien des valeurs par défaut, et \texttt{g} n'est pas virtuelle.}
\end{solution}

\begin{solution}{iq:iv3:cpp:35}
Il affiche \texttt{1 8 4 8 4}, \texttt{8 8 16} et \texttt{16 32 32 24}. Un objet doit avoir une adresse unique, si bien que deux éléments d'un tableau de \texttt{Empty} ont besoin d'adresses différentes, et la taille doit être d'au moins 1. \texttt{Base1} fait 4 et \texttt{Member} fait 8 grâce à l'optimisation de la base vide.

Une base vide peut partager une adresse avec le premier membre, car un sous-objet de base n'a pas besoin d'être distinct de cette façon. Un membre vide, en revanche, occupe un octet, et le remplissage arrondit \texttt{Member} à 8. L'attribut \texttt{[[no\_unique\_address]]} de C++20 accorde le même traitement au membre, ce qui explique que \texttt{Attr} fasse 4. \texttt{Virt} fait 8 pour son pointeur vers une table virtuelle.

Une lambda sans capture a un type vide, et son appel est connu à partir du type. Le \texttt{unique\_ptr} de libstdc++ stocke le deleter de façon qu'il n'occupe rien, si bien que la taille reste de 8. Un deleter sous forme de pointeur de fonction est une valeur réelle, choisie à l'exécution, et doit être stocké, ce qui donne 16.

Un \texttt{shared\_ptr} contient deux pointeurs~: un vers l'objet, et un vers le bloc de contrôle qui contient les compteurs. Deux sont nécessaires, car le pointeur vers l'objet peut différer de l'objet possédé, par exemple pour une classe de base ou avec un aliasing.

Un \texttt{std::function} fait 32 ici, puisqu'il contient un petit tampon pour l'objet appelable. Un \texttt{std::string} fait 32, avec une place pour 15 caractères, et un vecteur contient trois pointeurs. La norme ne fixe aucune de ces tailles.

\iqwrong{«~Une structure vide ne prend pas de place.~» Comme base, c'est en général le cas. Comme membre ou élément de tableau, elle en prend une, sauf indication contraire.}
\end{solution}

\begin{solution}{iq:iv3:cpp:36}
Il échoue à l'édition de liens, avec \texttt{undefined reference} vers \texttt{total<int>}. Chaque fichier est compilé seul. \texttt{main.cpp} ne voit que la déclaration, et suppose donc que quelqu'un d'autre fournit \texttt{total<int>}. \texttt{sum.cpp} voit le corps, mais rien dans ce fichier ne demande un \texttt{T}, si bien qu'il ne génère aucun code.

Un gabarit est une recette~: il ne produit du code que lorsqu'il est instancié avec un type, et seulement là où le compilateur voit à la fois la demande et la recette. Aucun fichier ici ne voit les deux.

Deux corrections. La plus courante consiste à placer le corps dans l'en-tête, afin que chaque appelant puisse l'instancier. Une définition de gabarit peut apparaître dans plusieurs fichiers, et l'éditeur de liens n'en conserve qu'un exemplaire.

La seconde est l'instanciation explicite. On laisse le corps dans \texttt{sum.cpp} et on ajoute à cet endroit la ligne \texttt{template int total(const std::vector<int>\&);}, qui fait générer cette seule version. Cela allège les en-têtes et accélère la compilation, et notre construction de test s'effectue ainsi, avec édition de liens.

Le coût, c'est la liste fermée. Un appelant qui utilise \texttt{std::vector<long>} obtient la même erreur d'édition de liens, comme le montre le test, tant que quelqu'un n'ajoute pas une seconde ligne dans \texttt{sum.cpp}. Un utilisateur d'une bibliothèque compilée ne peut pas en ajouter du tout.

\iqwrong{«~\texttt{export} règle le problème.~» Le sens de ce mot-clé pour les gabarits a été supprimé en C++11. Les modules de C++20 constituent la voie moderne, et ils exigent la prise en charge par le système de construction.}
\iqfollow{Quel est le coût de la forme à en-tête~?}{Le temps de compilation~: chaque fichier qui l'inclut analyse le corps et peut l'instancier, et \texttt{extern template} peut réduire ce coût.}
\end{solution}

\begin{solution}{iq:iv3:cpp:37}
Il affiche \texttt{integer real integer}, puis \texttt{10}. Chaque \texttt{kind} n'est viable que si son concept est vérifié pour le type de l'argument, et donc exactement une surcharge survit. \texttt{notional(1, 2)} échoue à l'appel avec une erreur indiquant que la contrainte \texttt{Priced<int>} n'est pas satisfaite, et quelle exigence a échoué.

Un concept est un prédicat nommé à la compilation portant sur des types. \texttt{requires(const T\& t) \{ \dots \}} teste si les expressions qu'il contient compileraient, et il peut aussi vérifier le type du résultat. Ainsi, \texttt{Priced} exige que \texttt{t.price()} soit valide et se convertisse en \texttt{double}. Ici, \texttt{int} n'a pas de membre \texttt{price}. Si deux surcharges contraintes s'appliquent toutes les deux, la plus contrainte est choisie.

Sans contrainte, le gabarit serait accepté pour \texttt{int}, et l'erreur surgirait dans son corps, à la ligne \texttt{t.price() * qty}. Dans une bibliothèque profonde, ce message peut s'étendre sur des centaines de lignes, et désigner un endroit très éloigné de l'erreur.

Avant C++20, le moyen courant était SFINAE~: un paramètre de gabarit supplémentaire, dont la valeur par défaut est \texttt{std::enable\_if\_t<std::is\_integral\_v<T>{}>}, retire la surcharge sans bruit lorsque la déduction échoue. Cela fonctionne, mais se lit mal, et deux surcharges qui ne diffèrent que par une telle valeur par défaut forment une redéfinition~; il faut donc un paramètre factice pour les distinguer. \texttt{if constexpr} couvre plus simplement certains cas, à partir de C++17.

\iqwrong{«~Un concept est une sorte de classe de base.~» Il n'a pas d'objets et n'ajoute aucun appel virtuel, et un type n'a besoin d'aucune déclaration pour le satisfaire.}
\end{solution}

\begin{solution}{iq:iv3:cpp:38}
Il affiche \texttt{300 2 24}, \texttt{1}, \texttt{bad access} et \texttt{0 1}. Un variant contient son alternative sur place, sans allocation~: 24~octets ici, pour la plus grande alternative de 16~octets alignée sur 8, plus l'indice. \texttt{get\_if} renvoie un pointeur nul pour le mauvais type, et \texttt{get} lève une exception.

Par rapport à une classe de base et trois classes dérivées, il n'y a ni tas ni pointeur à suivre, si bien qu'un vecteur de messages est contigu. \texttt{std::visit} fait vérifier par le compilateur que chaque type est traité, puisque \texttt{Overload} ne compile pas si l'un d'eux manque. Ajouter un type casse donc chaque visiteur tant qu'il n'est pas traité, ce qui est souhaitable ou constitue un coût, selon qui possède le code. Ajouter une opération est facile, puisqu'il s'agit d'un nouveau visiteur. Avec l'héritage, c'est l'inverse~: un nouveau type est facile, et une nouvelle opération touche chaque classe. Un ensemble fermé de messages connus convient au variant, et un ensemble ouvert, comme des stratégies en greffons, convient à l'héritage.

\texttt{std::visit} répartit selon l'indice, ce qui n'est pas nécessairement moins coûteux qu'un appel virtuel~; le gain à revendiquer porte donc sur la mémoire et la disposition. Nous n'avons rien chronométré ici.

Oui, il peut ne rien contenir. Changer de type détruit l'ancienne valeur avant de construire la nouvelle, et si cela lève une exception, il ne reste rien. C'est \texttt{valueless\_by\_exception()}, et \texttt{index()} renvoie \texttt{variant\_npos}, comme le montre le test lorsque \texttt{Fragile} possède un membre \texttt{std::string}. libstdc++ évite ce cas pour un petit type trivialement copiable, en construisant d'abord un temporaire~: avec un membre \texttt{int} à la place de la chaîne, le même programme affiche \texttt{0 0}.

\iqwrong{«~Un variant ne peut pas être vide.~» Il peut l'être, mais uniquement dans cet état d'échec, et \texttt{std::visit} sur lui lève une exception.}
\end{solution}

\begin{solution}{iq:iv3:cpp:39}
Il faut fournir une fonction de hachage et un test d'égalité~; sans fonction de hachage, la table ne compile pas. Le plus grand compartiment contient 50 entrées avec la fonction faible et, dans l'exécution présentée ici, 4 avec la fonction mélangée. La fonction faible donne des réponses correctes, mais coûte du temps~: chaque recherche parcourt une chaîne de jusqu'à 50 clés.

L'égalité est nécessaire, car des clés différentes peuvent partager un même hachage, et la table doit pouvoir les distinguer. C++20 peut la générer avec \texttt{= default}. Une fonction de hachage peut être passée en argument de gabarit, comme ici, ou sous forme de spécialisation de \texttt{std::hash<Key>}.

\texttt{Weak} n'utilise que le symbole, si bien que les 50 plateformes d'un symbole ont des hachages égaux et tombent dans un seul compartiment, quelle que soit la taille de la table. La table retrouve tout de même la bonne entrée en comparant les clés. Une fonction de hachage défaillante est donc un bogue de performance et jamais un bogue de correction, tant que des clés égales ont le même hachage. C'est l'unique règle stricte~: si deux clés égales avaient des hachages différents, les recherches échoueraient.

\texttt{Mixed} regroupe les deux champs dans 64 bits et multiplie par une constante impaire. libstdc++ prend le hachage modulo un nombre premier de compartiments, si bien que la valeur empaquetée seule suffirait à peu près~; mais une table réduite par un masque n'utiliserait que les bits de poids faible, et la multiplication n'y aide pas, car elle ne propage l'information que vers le haut~: les bits de poids faible du produit ne dépendent toujours que de la plateforme. Une telle table doit prendre les bits de poids fort du produit, ou faire d'abord un XOR des bits de poids fort vers le bas.

Le pire cas, à la limite, est adversarial~: si un attaquant peut choisir des clés ayant un même hachage, chaque opération devient linéaire.

\iqwrong{\texttt{h(symbol) \^{} h(venue)}. Le XOR est symétrique, et deux champs égaux s'annulent, si bien que (1, 2) et (2, 1) entrent en collision, et (5, 5) se hache en 0.}
\end{solution}

\begin{solution}{iq:iv3:cpp:40}
Il affiche \texttt{2 1}, \texttt{2 3}, \texttt{2} et \texttt{False True}. Le boxing consiste à copier une valeur de type valeur dans un nouvel objet sur le tas, afin qu'elle puisse être conservée par l'intermédiaire de \texttt{object} ou d'une interface. C\# le fait dans \texttt{object box = a}, dans \texttt{ICounter c = a}, et pour chacun de \texttt{o1} et \texttt{o2}.

Ainsi, \texttt{box} contient une copie avec \texttt{Qty} égal à 1, et la modification ultérieure de \texttt{a} ne l'atteint pas. \texttt{c.Add()} modifie la boîte, si bien que \texttt{a} reste à 2, tandis que la boîte affiche 3. Le piège tient à ce que la méthode semble modifier \texttt{a}. La lecture d'un élément de liste donne aussi une copie, si bien que la troisième ligne montre la liste inchangée. Et deux boîtes contenant 5 sont deux objets, ce qui donne la dernière ligne~: pas la même référence, mais des valeurs égales.

Chaque boîte est également une allocation que le ramasse-miettes (garbage collector) doit nettoyer, et c'est pourquoi un chemin critique (hot path) en C\# évite le boxing, en utilisant les génériques, où \texttt{List<Fill>} stocke les valeurs en ligne, et des appels contraints.

C++ n'a pas de boxing, car il n'a ni type de base commun ni objet implicite sur le tas. Un objet est une valeur, sauf demande contraire~; une copie est toujours visible comme copie, et la sémantique par référence est explicite~: \texttt{Fill\& ref} ou un pointeur. \texttt{list[0]} renvoie une référence, si bien qu'ici, contrairement à C\#, écrire à travers elle modifie l'élément. Lorsque C++ doit conserver des types différents derrière une même interface, il le dit explicitement avec un pointeur, un \texttt{std::any} ou un \texttt{std::variant}.

\iqwrong{«~Une \texttt{struct} C\# et une \texttt{struct} C++ sont la même idée.~» Une \texttt{struct} C++ ne diffère d'une \texttt{class} que par l'accès par défaut. En C\#, le mot-clé choisit la sémantique de valeur ou de référence.}
\end{solution}


\begin{solution}{iq:iv3:cpp:41}
Le \Cref{lst:iv3:cp:growth} affiche \texttt{1 2 4 8 16 32 64 128 256 512 1024}, puis \texttt{1023}. Pour 1000 insertions, le doublement fait 1023 déplacements, et l'ajout de 16 à chaque fois en fait 31\,248~; cet écart s'accroît avec $n$. \texttt{v.push\_back(v[0])} dans le vecteur naïf lit une mémoire déjà libérée, ce qu'AddressSanitizer signale comme \texttt{heap-use-after-free}.

Le doublement déplace $1 + 2 + \dots + 512 = 1023$ éléments, soit à peu près un par insertion, si bien que chaque insertion coûte en moyenne un nombre constant de déplacements. Un incrément fixe $k$ déplace $k + 2k + \dots$, soit environ $n^2 / 2k$ au total, ce qui est quadratique~: notre décompte est à peu près quadruplé de 1000 à 2000 insertions. L'incrément économise de la mémoire et fait perdre du temps.

Un facteur $f$ fait environ $n / (f - 1)$ déplacements~; ainsi, le facteur 1.5 en fait deux fois plus que 2, et un facteur proche de 1 devient le cas de l'incrément fixe. L'avantage d'un facteur plus petit est la réutilisation. Avec le doublement, chaque nouveau bloc est plus grand que tous les blocs libérés réunis. Avec 1.5, les blocs valent $1, 1.5, 2.25, 3.38, 5.06, 7.59$, et le sixième tient dans l'espace des quatre premiers, qui totalisent 8.13, si l'allocateur fusionne les blocs libres voisins. Sous le nombre d'or, cette réutilisation finit par devenir possible.

Dans le code naïf, \texttt{x} pointe dans l'ancien bloc, et \texttt{grow} supprime ce bloc avant que \texttt{x} ne soit lu. La version correcte construit donc d'abord le nouvel élément dans le nouveau bloc, et ne déplace ensuite les anciens et ne les libère qu'après.

\iqwrong{«~La croissance coûte $O(n)$ par insertion, puisque les éléments sont copiés.~» C'est le pire pas isolé. La moyenne sur toutes les insertions est constante.}
\iqfollow{Qu'y a-t-il d'autre de faux dans \texttt{new T[c]} ici~?}{Il exige un constructeur par défaut, et construit $c$ objets que personne n'a demandés.}
\end{solution}

\begin{solution}{iq:iv3:cpp:42}
Un membre \texttt{T value} devrait être construit à chaque fois que l'optional l'est, ce qui exige un constructeur par défaut et un objet que personne n'a demandé. Le stockage est donc un tableau d'octets bruts, de la taille et de l'alignement de \texttt{T}, et l'objet apparaît par placement \texttt{new}. Le \Cref{lst:iv3:cp:opt} ne lève l'indicateur qu'après la réussite du constructeur. Le programme mal aligné n'est pas correct, même s'il s'exécute.
\omcode{code/interviews-3/03-cpp/cpp/opt.hpp}{5}{39}{Un petit optional.}\label{lst:iv3:cp:opt}

\texttt{alignas(T)} compte, car un tableau de \texttt{unsigned char} a par lui-même un alignement de 1. Si l'indicateur était levé en premier et que le constructeur lève une exception, le destructeur détruirait un objet qui n'a jamais existé. Un test avec cet ordre échoue, tandis que l'ordre du listing réussit face à \texttt{std::optional} sur 20\,000 opérations aléatoires, sous AddressSanitizer, avec parmi elles un constructeur qui lève une exception et un type aligné sur 64 octets.

\texttt{std::launder} est nécessaire, car un transtypage de pointeur depuis le tableau d'octets ne pointe pas, à lui seul, vers l'objet créé à l'intérieur. L'affectation prend son argument par valeur, si bien que l'auto-affectation est sûre, puisque la copie existe avant \texttt{reset}.

Si \texttt{emplace} lève une exception, l'ancienne valeur a déjà disparu et l'optional est vide, ce qui est un état valide. C'est la garantie de base, et non la garantie forte.

Quant au \texttt{double} mal aligné~: UBSan signale \texttt{store to misaligned address} et \texttt{load of misaligned address}, et le programme affiche toujours 1.5. Un objet doit se trouver à une adresse convenablement alignée, donc ce comportement est indéfini. x86-64 tolère les déplacements simples, mais le compilateur peut vectoriser l'accès avec des instructions qui supposent l'alignement, et d'autres architectures peuvent provoquer un piège.

\iqwrong{«~\texttt{char buf[sizeof(T)]} suffit.~» C'est la bonne taille, mais au mauvais alignement.}
\end{solution}

\begin{solution}{iq:iv3:cpp:43}
Garder la liste libre (free list) à l'intérieur des emplacements libres eux-mêmes~: chacun est une union d'un pointeur suivant et des octets d'un \texttt{T}, si bien que la mémoire d'un emplacement libre contient le chaînage, et aucune mémoire supplémentaire n'est nécessaire. Le \Cref{lst:iv3:cp:pool} le fait. \texttt{create} retire la tête, et \texttt{destroy} la remet.
\omcode{code/interviews-3/03-cpp/cpp/pool.hpp}{6}{42}{Un pool à taille fixe avec une liste libre intrusive.}\label{lst:iv3:cp:pool}

Il y a un piège dans \texttt{create}. Le nouvel objet est construit par-dessus les octets qui contiennent \texttt{next}, si bien que \texttt{next} doit être lu avant que le constructeur ne s'exécute. Notre test construit la version qui le lit après~: le champ \texttt{id} écrase le chaînage avec 0, et le pool se déclare plein alors que 15 emplacements sont encore libres.

Si le constructeur lève une exception, rien n'a été créé, et l'emplacement retourne dans la liste~; une exception ne fait donc perdre aucune capacité. Le test lève une exception lors d'environ une création sur 8, et le compte des objets vivants ainsi que le test d'épuisement concordent toujours avec un modèle après 100\,000 opérations aléatoires.

Il n'appelle pas \texttt{operator new} dans cette exécution, et chaque emplacement est aligné pour \texttt{T}, ce que le test vérifie avec un type de 64 octets. Un emplacement fait au moins la taille d'un pointeur, quelle que soit la petitesse de \texttt{T}. L'emplacement libéré en dernier revient en premier, et c'est celui qui a le plus de chances de se trouver encore dans le cache.

Il ne détecte pas ces cas~: détruire deux fois un pointeur, ce qui corrompt la liste par un cycle~; utiliser un objet après \texttt{destroy}~; passer un pointeur provenant d'un autre pool. Et il est conçu pour un seul thread. Un second thread exige un verrou ou une liste sans verrou (lock-free), où le dépilement rencontre le problème ABA. Une version de débogage pourrait conserver un bitmap des emplacements vivants pour attraper le premier de ces cas.

\iqwrong{«~\texttt{T* p = new T[N]} et une pile de pointeurs.~» Cela construit les $N$ objets dès le départ et exige un constructeur par défaut, si bien que \texttt{create} n'aurait plus rien à construire.}
\end{solution}

\begin{solution}{iq:iv3:cpp:44}
Un allocateur n'a besoin que de \texttt{value\_type}, \texttt{allocate}, \texttt{deallocate}, de l'égalité, et d'un constructeur de conversion depuis le même allocateur pour un autre type. Le \Cref{lst:iv3:cp:arena} présente l'arène et cet allocateur. \texttt{deallocate} ne fait rien, car une arène à allocation linéaire libère tout d'un coup, en abandonnant le tampon.
\omcode{code/interviews-3/03-cpp/cpp/arena_alloc.hpp}{6}{35}{Une arène à allocation linéaire et un allocateur par-dessus.}\label{lst:iv3:cp:arena}

\texttt{std::align} avance le pointeur au-delà de tout remplissage, et échoue lorsque le reste du tampon est trop petit, auquel cas on lève une exception. Le test prend 1 octet puis un \texttt{double} aligné, et obtient une adresse alignée.

Le constructeur de conversion existe à cause de ce qu'un conteneur alloue. On nomme \texttt{pair<const int, double>}, mais un \texttt{std::map} alloue des nœuds d'arbre, un type privé que l'on ne voit jamais. Il rebinde votre allocateur à ce type de nœud et en construit un à partir du vôtre, ce que rend possible le constructeur. L'égalité indique si un espace de stockage alloué par l'un peut être libéré par l'autre~; ici, elle compare donc l'arène.

L'allocateur contient un pointeur vers l'arène et non l'arène elle-même, car les allocateurs sont copiés, et toutes les copies doivent puiser dans un seul tampon.

Les anciens blocs sont simplement perdus. Un vecteur qui grandit jusqu'à 100 entiers prend 1, 2, 4 jusqu'à 128, soit 255 entiers d'arène pour 100 éléments. Il faut donc appeler \texttt{reserve} lorsque l'on connaît la taille, ou utiliser un pool qui réutilise la mémoire.

Le test remplit un vecteur et une map à partir du tampon, ne fait aucun appel au \texttt{operator new} global, obtient les bonnes sommes, et voit \texttt{bad\_alloc} lorsqu'une requête est trop grande. Le tampon doit survivre à tous les conteneurs qui l'utilisent, ce que le type ne garantit pas.

\iqwrong{«~Les conteneurs appellent \texttt{allocate} pour les éléments que j'ai nommés, donc un seul type suffit.~» C'est vrai pour un vecteur, pas pour un conteneur à base de nœuds.}
\end{solution}

\begin{solution}{iq:iv3:cpp:45}
Il affiche \texttt{1020 400}, puis \texttt{1 0}, puis \texttt{bad\_alloc}. Sans \texttt{reserve}, le vecteur demande 1, 2, 4 jusqu'à 128 entiers, soit 255 entiers ou 1020 octets, et l'arène ne réutilise jamais rien. Avec \texttt{reserve(100)}, il ne demande qu'une fois, 400 octets.

La deuxième ligne montre que la ressource d'une chaîne est le tampon, et que le \texttt{new} global n'a été appelé aucune fois. Un \texttt{pmr::vector} qui construit un \texttt{pmr::string} transmet son propre allocateur, si bien que chaque élément utilise la même ressource. On parle de construction à allocateur propagé, et c'est la raison pour laquelle les conteneurs imbriqués sont si faciles à utiliser ici.

La dernière ligne vient de la ressource amont. Lorsque le tampon de 64 octets est épuisé, la ressource monotone demande davantage à son amont, et \texttt{null\_memory\_resource} lève toujours une exception. Cela transforme «~ce chemin ne doit pas atteindre le tas~» en un échec vérifiable. Avec l'amont par défaut, la requête irait discrètement sur le tas.

Le gain par rapport à un allocateur template se situe dans le type. \texttt{std::vector<int, ArenaAlloc<int>{}>} et \texttt{std::vector<int>} sont deux types différents, si bien qu'une fonction doit être un template pour accepter les deux. Tout \texttt{pmr::vector<int>} est un seul type, quelle que soit sa ressource, et le choix se fait donc à l'exécution.

Le coût est un pointeur dans chaque conteneur, et un appel virtuel à chaque allocation, ce qui pèse peu lorsque les allocations sont rares. Il existe aussi un piège. La construction par copie d'un \texttt{pmr::vector} n'hérite pas de la ressource de la source~: la copie utilise la ressource par défaut, ce que montre notre test, si bien qu'une copie peut atteindre le tas sans que personne ne s'en aperçoive.

\iqwrong{«~Le vecteur libère ses anciens blocs dans le tampon.~» \texttt{do\_deallocate} sur une ressource monotone ne fait rien, et la mémoire ne revient qu'à la libération ou à la destruction de la ressource.}
\end{solution}

\begin{solution}{iq:iv3:cpp:46}
\texttt{a = a} efface silencieusement les données, et le programme affiche 0 au lieu de 7. \texttt{a = b} avec un \texttt{get} qui lève une exception laisse \texttt{d\_} pointant vers une mémoire libérée, si bien que le destructeur la libère de nouveau, et AddressSanitizer signale \texttt{attempting double-free}. Le \Cref{lst:iv3:cp:assigngood} est correct dans les deux cas, grâce à la copie puis à l'échange.
\enlargethispage{-3\baselineskip}%
\omcode{code/interviews-3/03-cpp/cpp/assign_good.cpp}{12}{30}{Copie, puis échange.}\label{lst:iv3:cp:assigngood}

Lors de l'auto-affectation, \texttt{o} et \texttt{*this} sont un seul objet. \texttt{delete[] d\_} libère les données, puis \texttt{d\_ = get(n\_)} installe un bloc neuf mis à zéro, et \texttt{std::copy} recopie ce bloc sur lui-même. Aucune mémoire libérée n'est lue, si bien qu'AddressSanitizer reste silencieux, et les données ont simplement disparu. Cette auto-copie viole aussi la précondition de \texttt{std::copy}~; formellement, elle est donc indéfinie~; avec cette bibliothèque, c'est un \texttt{memmove} sans effet néfaste. Un \texttt{if (this != \&o)} corrigerait cette faute et laisserait l'autre en place.

La seconde faute tient à l'ordre~: l'ancien état est détruit avant que l'on sache que le nouveau existe. Si \texttt{get} lève une exception, \texttt{d\_} contient encore le pointeur supprimé.

Dans la nouvelle version, la copie est faite en premier, alors que \texttt{*this} n'est pas touché, et c'est la seule ligne qui puisse lever une exception. L'échange ne le peut pas, et le destructeur de \texttt{copy} libère les anciennes données. L'auto-affectation n'a besoin d'aucun contrôle. Le test affiche \texttt{7}, \texttt{caught 7} et \texttt{8}~: la valeur survit à l'auto-affectation, et une affectation échouée la laisse inchangée.

Le coût tient à une allocation à chaque affectation, même lorsque le bloc existant est assez grand. Un \texttt{std::vector} l'évite en copiant dans sa capacité existante, et renonce ainsi à la garantie forte pour son affectation par copie~: si l'affectation d'un élément lève une exception en cours de route, le vecteur contient un mélange. Sur un chemin critique, c'est le bon compromis, à condition que l'appelant puisse s'en accommoder.

\iqwrong{«~Ajoutez \texttt{if (this == \&o) return *this;} et c'est corrigé.~» Cela traite un cas qui arrive à peine, et laisse en place la défaillance qui, elle, arrive.}
\end{solution}

\begin{solution}{iq:iv3:cpp:47}
Il affiche \texttt{1}, puis \texttt{3f800000}. Un \texttt{int*} et un \texttt{float*} sont supposés ne pas pointer vers le même objet, si bien que le compilateur renvoie 1 sans relire \texttt{*a}. Avec \texttt{-fno-strict-aliasing}, il ne fait pas cette supposition, relit la valeur, et la première ligne devient 0.

La règle~: un objet ne peut être lu qu'à travers son propre type, une version signée ou non signée de celui-ci, ou un \texttt{char}, un \texttt{unsigned char} ou un \texttt{std::byte}. Lire un \texttt{int} à travers un \texttt{float} est indéfini~; c'est donc \texttt{main} qui est fautif ici, et non l'optimiseur. Des hypothèses de ce genre permettent au compilateur de garder des valeurs dans des registres à travers des écritures sans rapport, si bien qu'une boucle qui écrit à travers un \texttt{float*} n'a pas à relire un \texttt{int} à chaque tour.

La deuxième ligne montre la manière légale de voir les bits d'un flottant~: \texttt{std::bit\_cast} donne une nouvelle valeur aux mêmes bits, ici 0x3f800000 pour 1.0.

Pour un paquet, \texttt{*reinterpret\_cast<const std::uint32\_t*>(p)} enfreint deux règles~: l'aliasing et l'alignement, puisqu'un champ d'un paquet peut se trouver à n'importe quel décalage. La bonne façon consiste à déclarer un \texttt{std::uint32\_t v} puis \texttt{std::memcpy(\&v, p, 4)}. Cela semble lent, et ne coûte rien ici~: avec \texttt{-O2}, notre test trouve une seule instruction \texttt{movl (\%rdi), \%eax}, et aucun appel. Il faut inverser l'ordre des octets ensuite, si le format sur le réseau l'exige.

De nombreuses grandes bases de code compilent avec \texttt{-fno-strict-aliasing} comme filet de sécurité~; c'est un choix, et il rend le programme non portable.

\iqwrong{«~Une \texttt{union} des deux est la voie légale.~» C'est permis en C, et c'est indéfini en C++, bien que g++ le documente comme une extension.}
\iqfollow{Et plaquer une struct sur le tampon~?}{Les mêmes deux règles sont enfreintes, et la disposition en mémoire pose aussi la question du remplissage. Il faut copier les champs, ou faire un \texttt{memcpy} vers une struct dont la disposition a été vérifiée.}
\end{solution}

\begin{solution}{iq:iv3:cpp:48}
\texttt{read\_bad} se compile en un chargement et un retour, et ne contient aucun contrôle. \texttt{read\_good} teste le pointeur, et comporte un branchement qui renvoie $-1$. Le déréférencement de la première ligne indique au compilateur que \texttt{p} n'est pas nul, car déréférencer un pointeur nul est indéfini~; \texttt{if (!p)} ne peut donc jamais être vrai, et il est supprimé.

Le compilateur ne fait aucun effort particulier ici. Il propage un fait qu'il a appris, comme il le fait partout, puis élimine le code mort. Chaque étape est valide, et le résultat ne choque que parce que le contrôle était le filet de sécurité de l'auteur. Aucun avertissement n'est promis, et notre compilation avec \texttt{-Wall -Wextra} n'en produit aucun.

Avec \texttt{read\_bad(nullptr)} sur cette plateforme, le chargement lit l'adresse 0, et le processus s'arrête sur une faute de segmentation. Le contrôle était donc inutile même avant d'être supprimé, puisque le plantage survient à la ligne précédente. Toutefois, rien ne le garantit non plus~: si \texttt{v} n'était pas utilisé, le chargement pourrait disparaître lui aussi, et le programme poursuivrait dans un état que son auteur croyait impossible.

\texttt{-fno-delete-null-pointer-checks} conserve le contrôle, et notre test retrouve l'instruction \texttt{test} dans l'assembleur. Sur un système où l'adresse 0 peut être mappée, cela compte. Comme correctif pour ce code, c'est la mauvaise réponse~: c'est l'ordre des lignes qui est en cause, et l'option ne fait que le masquer pour un seul compilateur. Il faut placer le contrôle avant l'utilisation, comme le fait \texttt{read\_good}.

La leçon générale est que le comportement indéfini agit en amont comme en aval. Une opération future peut modifier ce que le compilateur suppose du passé~; «~cela ne plante que si on y arrive~» n'est donc pas fiable.

\iqwrong{«~L'optimiseur a supprimé un contrôle, c'est donc un bogue du compilateur.~» La norme lui permet de supposer l'absence de comportement indéfini.}
\end{solution}

\begin{solution}{iq:iv3:cpp:49}
Sans argument, AddressSanitizer signale \texttt{stack-use-after-scope} à la ligne 6. Avec trois arguments, le programme s'exécute sans erreur et affiche 4. \texttt{std::max} renvoie une référence vers l'argument qui l'emporte. Avec \texttt{argc} égal à 1, c'est le temporaire contenant 2, qui n'existe plus à la fin de la ligne 5. Avec \texttt{argc} égal à 4, c'est \texttt{argc} lui-même, qui est toujours vivant.

Le bogue dépend donc des données. Un test dont \texttt{argc} est supérieur à 2 passe, et rien ne signale que l'autre cas pend dans le vide.

La liaison de référence \texttt{m} n'étend aucune durée de vie. Le temporaire est lié au paramètre de la fonction, et il vit jusqu'à la fin de l'expression complète. Ce que renvoie la fonction est une nouvelle référence, et le compilateur ne sait pas qu'elle peut désigner ce temporaire.

\texttt{std::max} renvoie une référence pour ne pas imposer de copie. Elle convient à un type qui ne peut pas être copié ou dont la copie coûte cher, et son résultat peut être utilisé sur place, dans la même expression. C'est sûr tant que le résultat n'est pas conservé.

\texttt{std::max(1, 2.0)} ne compile pas, car \texttt{T} est déduit des deux arguments et vaut à la fois \texttt{int} et \texttt{double}. Aucune conversion n'est envisagée lors de la déduction. Il faut écrire \texttt{std::max<double>(1, 2.0)}.

La correction consiste ici à conserver une valeur~: \texttt{int m = std::max(argc, 2);}. Cela copie pendant que le temporaire est vivant, et pour un \texttt{int} la copie ne coûte rien. Il en va de même pour \texttt{std::min} et \texttt{std::clamp}.

\iqwrong{«~Une référence \texttt{const} prolonge le temporaire, c'est donc sûr.~» L'extension ne s'applique qu'à un temporaire lié directement par cette référence.}
\iqfollow{Et \texttt{std::max(a, b).size()} sur deux chaînes~?}{Cela va~: l'appel et l'utilisation se trouvent dans une même expression complète, donc rien n'est encore détruit.}
\end{solution}

\begin{solution}{iq:iv3:cpp:50}
Le mode \texttt{A} est une course aux données, et ThreadSanitizer la signale. Le compteur du bloc de contrôle est atomique, mais l'objet \texttt{shared\_ptr} ne l'est pas~: il se compose de deux pointeurs, et un thread les écrit pendant qu'un autre les lit. Le mode \texttt{B} passe par \texttt{std::atomic\_store} et \texttt{std::atomic\_load}, et il est propre.

Ce qui peut mal tourner est concret. Le lecteur charge le pointeur vers le bloc de contrôle, puis est mis en pause. L'écrivain remplace \texttt{config}, l'ancien compteur atteint zéro, et le bloc est libéré. Le lecteur se réveille, incrémente un compteur dans une mémoire libérée, et détient un objet qui n'existe plus. Il peut aussi voir le nouveau pointeur d'objet avec l'ancien bloc, puisque les deux ne sont pas lus ensemble.

La règle est la règle habituelle pour un type~: les lectures concurrentes d'un même objet sont sans danger, de même que tout accès à des objets distincts, et une écriture sur un objet accompagnée de toute autre opération sur ce même objet est une course. Le compteur atomique rend sûrs des objets distincts, même lorsqu'ils partagent un objet possédé. Quatre threads disposant chacun de leur copie ne constituent donc pas une course, et ThreadSanitizer n'en trouve aucune dans le \cref{lst:iv3:cp:sharedcopies}.
\enlargethispage{-6\baselineskip}%
\omcode{code/interviews-3/03-cpp/cpp/shared_copies.cpp}{6}{18}{Chaque thread dispose de sa propre copie.}\label{lst:iv3:cp:sharedcopies}

\texttt{std::atomic<std::shared\_ptr<int>{}>} est la voie de C++20, mais libstdc++ 11 ne la fournit pas~: g++ 11.4 s'arrête sur un \texttt{static\_assert}, car \texttt{shared\_ptr} n'est pas trivialement copiable. Elle est arrivée avec GCC 12. Nous avons donc utilisé les anciennes fonctions libres, que libstdc++ implémente en prenant l'un d'un petit ensemble de mutex, comme le montre son en-tête. Ces fonctions sont dépréciées en C++20.

Ce chemin prend donc un verrou, et même une simple copie d'un \texttt{shared\_ptr} est une opération de lecture-modification-écriture atomique sur une ligne que tous les lecteurs partagent. Un lecteur sollicité intensivement préfère garder sa propre copie, et la rafraîchir de temps en temps.

\iqwrong{«~\texttt{shared\_ptr} est sûr vis-à-vis des threads.~» Son compteur l'est. L'objet qui contient les deux pointeurs ne l'est pas.}
\end{solution}

\begin{solution}{iq:iv3:cpp:51}
\Cref{lst:iv3:cp:bq} en donne une réponse. Il y a deux variables de condition, de sorte qu'un push ne réveille que les consommateurs et un pop que les producteurs~; chaque attente porte un prédicat, car une attente peut se terminer alors que la condition n'est pas remplie. \texttt{close} positionne le drapeau sous le verrou, puis réveille tout le monde sur les deux.
\omcode{code/interviews-3/03-cpp/cpp/bq.hpp}{9}{44}{Une file bornée avec close.}\label{lst:iv3:cp:bq}

Avec une seule variable de condition, un push pourrait réveiller un autre producteur, qui trouve la file pleine et se rendort~; et si \texttt{notify\_one} choisit celui-là, le réveil est perdu, et le consommateur en attente se rendort. Employer \texttt{notify\_all} corrigerait cela, au prix de nombreux réveils inutiles.

Une attente peut rendre la main pour deux raisons sans que la condition soit vraie~: un réveil fallacieux, que la norme autorise, et un réveil volé, où un autre thread a pris l'élément avant que ce thread obtienne le verrou. La forme avec prédicat boucle et vérifie à nouveau.

Positionner le drapeau sous le verrou est essentiel. Sinon, un thread en attente pourrait vérifier son prédicat, ne voir aucune fermeture, et s'endormir juste après la notification, sans jamais se réveiller. La notification elle-même peut venir après le déverrouillage, de sorte qu'un thread réveillé ne se bloque pas aussitôt sur le mutex que l'auteur de la notification tient encore. On note que \texttt{pop} continue de renvoyer des éléments après \texttt{close}~: la file se vide, et ne renvoie rien que lorsqu'elle est vide.

Le test fait tourner 3 producteurs et 3 consommateurs sur une file de capacité 4. Il vérifie le nombre, 9000, et la somme exacte, 27\,009\,000, ce qui détecterait un élément perdu ou dupliqué. Le test s'exécute sous ThreadSanitizer. Il bloque ensuite un producteur sur une file pleine, et confirme que \texttt{close} le réveille avec \texttt{false}, et que l'élément déjà présent peut encore être extrait. Un blocage ferait échouer le test à son délai d'expiration.

\iqwrong{«~\texttt{if (q\_.empty()) wait}~» au lieu d'une boucle. La fonction rend la main sur une file vide dès le premier réveil volé.}
\end{solution}

\begin{solution}{iq:iv3:cpp:52}
\Cref{lst:iv3:cp:choose} construit le résultat un facteur à la fois, en simplifiant les facteurs communs avant de multiplier. Il concorde avec un calcul exact pour tous les $n \le 75$, signale un dépassement exactement sur les cas qui ne tiennent pas, et le premier dépassement apparaît à $n = 68$, pour $k = 31$.
\omcode{code/interviews-3/03-cpp/cpp/choose.cpp}{1}{17}{Un coefficient binomial qui signale le dépassement.}\label{lst:iv3:cp:choose}

La formule factorielle échoue immédiatement, puisque $21!$ dépasse $2^{64}$, alors que $\binom{21}{10}$ vaut seulement 352\,716.

La boucle conserve $r = \binom{n - k + i}{i}$ après l'étape $i$. Pour atteindre la valeur suivante, il faut calculer $r \times (n - k + i) / i$, qui est une division exacte. Multiplier d'abord peut provoquer un dépassement alors que le résultat tient. On pose donc $g = \gcd(r, i)$. Alors $r / g$ et $i / g$ n'ont pas de facteur commun, de sorte que $i / g$ doit diviser $n - k + i$ exactement. On obtient $(r / g) \times ((n - k + i) / (i / g))$, produit de deux entiers égal au binomial suivant, sans qu'aucune valeur plus grande n'intervienne en cours de route.

Pourquoi aucun dépassement n'est-il signalé à tort~? Ce produit vaut exactement $\binom{n - k + i}{i}$, et ces valeurs croissent avec $i$ jusqu'à $k$, puisque $k \le n / 2$ après prise du minimum. Donc si une étape déborde, alors $\binom{n}{k}$ déborde aussi, c'est la dernière d'entre elles.

Le plus petit $n$ est 68, puisque $\binom{67}{33}$ vaut environ $1.42 \times 10^{19}$ et tient, et que $\binom{68}{31}$ vaut environ $2.19 \times 10^{19}$, au-dessus de $2^{64} \approx 1.84 \times 10^{19}$. Le test vérifie aussi les cas limites~: $k > n$ donne 0, $k = 0$ donne 1, et $n = 2^{64} - 1$ avec $k = 1$ renvoie $n$, tandis qu'avec $k = 2$ il y a dépassement.

\texttt{\_\_builtin\_mul\_overflow} est une extension de g++ et de Clang. Une forme portable teste d'abord \texttt{a > max / b}.

\iqwrong{\texttt{r = r * (n - k + i) / i} «~puisque la division est exacte~». La division est exacte, mais la multiplication qui la précède peut déborder alors que le résultat juste tient.}
\end{solution}

\begin{solution}{iq:iv3:cpp:53}
Avec \texttt{noexcept} et \texttt{destructor}, le programme affiche \texttt{terminate called} suivi du message, et se termine par \texttt{SIGABRT}. Avec \texttt{plain}, il affiche \texttt{caught plain}. Le bloc \texttt{try} ne reçoit jamais les deux premiers, car l'exception ne peut pas sortir de la fonction qui a promis de ne rien lever.

\texttt{noexcept} est une promesse vérifiée à la frontière de la fonction. Lorsque la recherche d'un gestionnaire atteint la frontière d'une fonction qui ne lève pas d'exception, le langage appelle \texttt{std::terminate} au lieu de poursuivre vers l'extérieur. La norme ne tranche pas la question de savoir si la pile est déroulée au préalable, de sorte que les destructeurs peuvent ne pas avoir été exécutés.

Les destructeurs sont \texttt{noexcept} par défaut depuis C++11, sans qu'on l'écrive, sauf si le destructeur d'un membre ou d'une base est déclaré pouvant lever une exception. Ainsi \texttt{\textasciitilde{}Closer()} ne lève rien, tout comme \texttt{promised}. Le compilateur le voit et émet un avertissement avec \texttt{-Wterminate}, activé par défaut. Le listing désactive cet avertissement sur une ligne non montrée, pour illustrer le propos.

On peut l'autoriser~: \texttt{\textasciitilde{}Closer() noexcept(false)}. Notre test le compile, et il affiche \texttt{caught from destructor}. Mais si un destructeur lève une exception alors qu'une autre exception déroule déjà la pile, il y en aurait deux à la fois, et le programme se termine, sans moyen d'en attraper une. Le test le montre aussi. Un destructeur susceptible de s'exécuter pendant un déroulement ne doit donc pas lever d'exception.

Un objet qui ferme un fichier ou un socket doit signaler l'échec autrement. Il faut proposer une méthode \texttt{close()} qui peut lever une exception, ou qui renvoie un statut pour l'appelant qui s'en soucie, et laisser le destructeur avaler et journaliser ce qui reste.

\iqwrong{«~Envelopper le code dans \texttt{try} et \texttt{catch (.{}.{}.)} dans \texttt{main}.~» C'est le code montré. Le contrôle n'y revient jamais.}
\iqfollow{Est-ce que \texttt{noexcept} rend le code plus rapide~?}{Il supprime un peu de travail pour l'appelant, et il permet aux conteneurs de déplacer les éléments. Le second point est celui où il compte le plus.}
\end{solution}

\begin{solution}{iq:iv3:cpp:54}
Il affiche 62~: \texttt{static\_call} donne $(10 + 1) \times 2 = 22$, et les deux appels virtuels donnent 22 et 18. \texttt{Momo} dérive de \texttt{Strat<Momo>}, donc la base connaît son type dérivé à la compilation, et \texttt{static\_cast<Derived*>(this)} est un appel simple et résolu. L'assembleur de \texttt{static\_call} ne contient aucun \texttt{call}, alors que \texttt{virtual\_call} contient \texttt{call *(\%rax)}, un appel indirect par la table virtuelle.

C'est le curiously recurring template pattern. Une classe de base gabarit prend la classe dérivée comme paramètre. Le transtypage est sûr si l'objet est bien un \texttt{Momo}, ce qui tient tant que seul \texttt{Momo} dérive de \texttt{Strat<Momo>}.

Le gain vient surtout de l'inlining. Le compilateur voit quelle fonction est appelée, il en recopie donc le corps, ce qui ouvre d'autres optimisations à travers la frontière de l'appel. L'appel indirect lui-même est généralement peu coûteux lorsqu'il est bien prédit. Ainsi la différence dépend de ce que fait la fonction appelée~: pour une petite fonction dans une boucle serrée elle peut être grande, et pour une fonction qui fait un vrai travail elle est faible. Nous n'avons rien mesuré ici, donc il faut mesurer dans son propre code.

Sans \texttt{signal}, l'erreur apparaît à la ligne de \texttt{on\_tick}, et seulement si \texttt{on\_tick} est utilisée, puisqu'un membre d'une classe gabarit est instancié à l'usage. Un concept pourrait donner un message plus clair.

\texttt{Strat<Momo>} et \texttt{Strat<Other>} sont des types sans lien, car il n'existe pas de base commune que le conteneur pourrait contenir. Le CRTP convient donc à une stratégie connue à la compilation. Pour choisir parmi les stratégies à l'exécution, il faut utiliser une fonction virtuelle, ou un \texttt{std::variant} et \texttt{std::visit}. Un mélange courant est un appel virtuel par message, avec du code à répartition statique à l'intérieur.

\iqwrong{«~Le CRTP est toujours plus rapide, donc supprimer toutes les fonctions virtuelles.~» Un compilateur peut aussi supprimer la répartition virtuelle lorsqu'il connaît le type, comme avec \texttt{final}, et un appel virtuel bien prédit coûte peu.}
\end{solution}

\begin{solution}{iq:iv3:cpp:55}
Il affiche \texttt{42 2 42 1000}. Seule \texttt{int d = twice(argc);} des quatre se compile. Une fonction \texttt{constexpr} peut s'exécuter à la compilation ou à l'exécution, et \texttt{consteval} doit s'exécuter à la compilation. Une variable \texttt{constexpr} doit avoir un initialiseur constant.

Ainsi \texttt{twice(21)} servant à initialiser une variable \texttt{constexpr} s'exécute dans le compilateur, et \texttt{twice(argc)} s'exécute dans le programme, puisque \texttt{d} n'exige aucune constante. \texttt{always(argc)} échoue, car un appel \texttt{consteval} doit donner une constante, et \texttt{argc} n'en est pas une. \texttt{constexpr int f = twice(argc)} échoue, car la variable exige ce que l'argument ne peut pas fournir.

Le dernier est le plus instructif. Doubler 2\,000\,000\,000 déborde \texttt{int}. À l'exécution, cela serait un comportement indéfini, compilé en silence. Dans une expression constante, la norme exclut toute opération indéfinie, et le compilateur doit donc la diagnostiquer. L'évaluation constante est ainsi un contrôle gratuit du comportement indéfini, pour les valeurs que l'on peut calculer à la compilation, ce qui est une des raisons d'écrire un \texttt{static\_assert} sur une table.

La table \texttt{powers} est construite par une boucle à l'intérieur du compilateur, et elle entre dans le programme comme donnée en lecture seule. Aucun code d'initialisation n'est nécessaire, il n'y a pas de problème d'ordre d'initialisation, et \texttt{powers[3]} se contente de lire une constante. Le \texttt{static\_assert} qui vérifie que \texttt{powers[18]} vaut $10^{18}$ contrôle la table à chaque construction. Notons la garde dans la boucle, qui l'empêche de calculer $10^{19}$~: sans elle, la dernière multiplication déborde, et la construction échoue.

Le coût se situe à la compilation, et dans les données du binaire. Une table utilisée dans un chemin critique peut aussi entrer en concurrence pour le cache, ce qu'un petit calcul ne fait pas, de sorte qu'une table n'est pas toujours le choix le plus rapide.

\iqwrong{«~\texttt{constexpr} sur une fonction signifie qu'elle s'exécute à la compilation.~» Cela ne fait que le permettre. C'est l'appel qui décide.}
\iqfollow{Et \texttt{constinit}~?}{Il exige une initialisation constante d'une variable statique, et la variable peut ensuite changer, contrairement à \texttt{constexpr}.}
\end{solution}

\begin{solution}{iq:iv3:cpp:56}
\Cref{lst:iv3:cp:vec1,lst:iv3:cp:vec2} donnent une implémentation. Une réallocation procède ainsi~: allouer le nouveau bloc, construire le nouvel élément à sa place finale, transférer les anciens éléments, les détruire puis libérer l'ancien bloc, et seulement alors mettre à jour les membres. Le vecteur se déplace lorsque le déplacement ne peut pas lever d'exception ou lorsqu'il n'existe pas de copie, et il copie sinon.
\omcode{code/interviews-3/03-cpp/cpp/vec.hpp}{8}{39}{Un vecteur~: stockage, copie, déplacement.}\label{lst:iv3:cp:vec1}
\omcode{code/interviews-3/03-cpp/cpp/vec.hpp}{40}{69}{Un vecteur~: croissance et accès.}\label{lst:iv3:cp:vec2}

\texttt{new T[c]} exigerait un constructeur par défaut, et construirait $c$ objets, laissant une capacité avec des objets vivants. Ici la mémoire provient de \texttt{std::allocator} sous forme de stockage brut, les objets étant créés par placement \texttt{new} et détruits à la main. La taille \texttt{n\_} indique exactement quels emplacements contiennent des objets, ce qui constitue l'invariant sur lequel repose la classe.

L'ordre satisfait les deux exigences. Le nouvel élément vient en premier, pendant que l'ancien bloc est encore vivant, de sorte qu'un argument qui renvoie à l'intérieur du vecteur, \texttt{v[0]}, reste valide. Notre test le fait 40 fois avec une chaîne de 40 caractères. Et jusqu'aux dernières étapes, rien de visible n'a changé. Si le constructeur du nouvel élément lève une exception, ou si une copie pendant le transfert en lève une, les blocs catch défont ce qui a été fait dans le nouveau bloc, et le vecteur reste tel qu'il était. \texttt{uninitialized\_copy\_n} détruit son propre travail partiel lorsqu'il échoue.

Un déplacement qui lève une exception laisserait certains anciens éléments déplacés, de sorte que le transfert copie lorsqu'il le peut, comme le fait le vrai vecteur. Un type qui ne supporte que le déplacement n'a pas le choix, et le test contient des \texttt{unique\_ptr}.

Pour tester la garantie, il faut utiliser un type dont la copie lève une exception à un nombre choisi d'appels, à chaque position tour à tour, puis comparer le contenu et la capacité après coup avec une copie prise avant. Vérifier ensuite le nombre d'objets vivants, et exécuter sous AddressSanitizer. Ici il exécute 20\,000 étapes aléatoires contre \texttt{std::vector}. Une version qui force le déplacement échoue toujours, ce que le test confirme.

\iqwrong{\texttt{put(); d\_ = fresh;} avant le transfert, ou le nouvel élément construit en dernier. L'un ou l'autre casse \texttt{v.emplace\_back(v[0])} ou la garantie.}
\end{solution}

\begin{solution}{iq:iv3:cpp:57}
L'incrémentation peut être relâchée. La décrémentation exige \texttt{acq\_rel}. Le \Cref{lst:iv3:cp:rc} est la classe, avec le compteur et l'objet dans un même bloc, et elle concorde avec \texttt{shared\_ptr} sur 20\,000 opérations aléatoires. Elle est propre sous AddressSanitizer et ThreadSanitizer, avec quatre threads qui se disputent la suppression de la dernière référence.
\omcode{code/interviews-3/03-cpp/cpp/rc.hpp}{5}{33}{Un pointeur à comptage de références.}\label{lst:iv3:cp:rc}

L'incrémentation est relâchée, car une nouvelle référence ne peut être créée qu'à partir d'une référence existante, ce qui maintient l'objet en vie pendant toute la durée, et l'incrémentation n'a pas à ordonner quoi que ce soit d'autre.

La décrémentation a deux rôles. Un thread qui a utilisé l'objet et abandonne sa référence doit avoir terminé tous ses accès avant que le compteur ne baisse, ce qui exige \texttt{release}. Et le thread qui atteint zéro et supprime l'objet doit voir tous ces accès comme achevés, ce qui exige \texttt{acquire}. Si elle est relâchée, la dernière lecture d'un autre thread peut être ordonnée après le \texttt{delete}, et le destructeur peut voir des membres périmés. Un raffinement courant utilise \texttt{release} sur la décrémentation, avec une barrière acquire seulement lorsque le compteur atteint zéro, ce qui évite l'acquire à chaque autre décrémentation.

Sur x86-64, les opérations atomiques de lecture-modification-écriture sont des barrières complètes, de sorte qu'une version relâchée fonctionne normalement, et un test de résistance n'échouera pas sur le matériel. C'est pourquoi le test utilise ThreadSanitizer, qui raisonne sur le modèle mémoire et non sur la machine. Il compile la classe avec la décrémentation passée en relâchée, et signale une \texttt{data race} sur les données de l'objet. Le même pilote avec \texttt{acq\_rel} n'en signale aucune. Le test peut donc mettre en défaut le modèle. Nous n'avons pas exécuté cela sur une machine à ordonnancement faible.

Cela ne rend pas un objet \texttt{Rc} unique sûr à modifier depuis deux threads, car il y a alors une course aux données sur \texttt{b\_}, et \texttt{count()} peut être périmé au moment où il rend sa valeur. Cela ne protège pas non plus le \texttt{T}, dont les membres ont besoin de leur propre synchronisation. Il n'a ni compteur faible, ni suppresseur personnalisé, ni aliasing, et c'est ce qui le fait tenir en 8 octets.

\iqwrong{«~Chaque opération devrait être \texttt{seq\_cst}, par précaution.~» C'est correct, mais cela coûte sur les machines à ordonnancement faible, et montre qu'on n'a pas identifié ce que chaque opération exige.}
\end{solution}

\begin{solution}{iq:iv3:cpp:58}
\Cref{lst:iv3:cp:fn1,lst:iv3:cp:fn2} donnent une implémentation. On efface le type en écrivant des fonctions tant qu'on le connaît encore~: \texttt{Model<F, Small>} génère une table de pointeurs de fonctions, une table par type appelable, et \texttt{Fn} conserve un pointeur vers la table. Au-delà de l'appel, il faut \texttt{move\_to} et \texttt{destroy}.
\omcode{code/interviews-3/03-cpp/cpp/fn.hpp}{7}{35}{Un support d'appelable~: la table pour chaque type.}\label{lst:iv3:cp:fn1}
\enlargethispage{-3\baselineskip}%
\omcode{code/interviews-3/03-cpp/cpp/fn.hpp}{37}{63}{Un support d'appelable~: le côté public.}\label{lst:iv3:cp:fn2}

À l'intérieur du constructeur, \texttt{D} est connu, et les lambdas de \texttt{Model<D, small>} sont écrites avec lui. Ensuite, \texttt{Fn} ne voit que \texttt{void*} et trois pointeurs de fonctions. C'est l'effacement de type, une table virtuelle écrite à la main.

Le tampon ne contient un appelable que s'il est assez petit, non suraligné, et qu'il possède un déplacement \texttt{noexcept}. Cette dernière condition découle du déplacement de \texttt{Fn}, qui est lui-même \texttt{noexcept}~: un appelable dans le tampon doit réellement se déplacer lorsque \texttt{Fn} se déplace, alors qu'un appelable sur le tas n'a besoin que de la copie de son pointeur. Si un déplacement dans le tampon pouvait lever une exception, le déplacement de \texttt{Fn} ne pourrait pas honnêtement promettre le contraire. La promesse compte~: un vecteur en croissance déplace un type qu'il ne peut pas copier de toute façon, mais si un tel déplacement lève une exception à mi-chemin, les éléments déjà déplacés le restent, et la garantie forte du vecteur disparaît.

Une lambda à quatre captures de type \texttt{long} occupe 32 octets, et le test compte 0 allocation pour elle. Une lambda avec un tableau de huit éléments en occupe 64 et coûte 1 allocation. Une fois déplacée, la source teste à faux, et l'appelable déplacé fonctionne toujours, sans copier les 64 octets. Une lambda qui contient un \texttt{unique\_ptr} fonctionne, puisqu'on ne demande jamais de copie. Son test différentiel contre \texttt{std::function} exécute 20\,000 étapes aléatoires avec un compteur d'objets vivants, qui vaut 0 à la fin, sous AddressSanitizer.

\texttt{std::function} possède la copie, ce qui explique qu'il ne puisse pas contenir une lambda avec \texttt{unique\_ptr}. Il lève aussi \texttt{bad\_function\_call} lorsqu'il est vide, alors que l'appel de notre \texttt{Fn} vide est un comportement indéfini, et il prend en charge \texttt{target}. C++23 ajoute \texttt{std::move\_only\_function}, qui est le plus proche de cette classe.

\iqwrong{«~Utilisez une classe de base virtuelle et \texttt{new} pour chaque appelable.~» Cela fonctionne, et cela alloue à chaque fois. La table écrite à la main est ce qui permet le petit tampon.}
\end{solution}

\begin{solution}{iq:iv3:cpp:59}
Le \Cref{lst:iv3:cp:var} donne une implémentation. Le stockage a la taille la plus grande parmi les types, aligné sur le plus strict d'entre eux, ce que \texttt{alignas(Ts.{}.{}.)} et \texttt{std::max(\{sizeof(Ts).{}.{}.\})} expriment directement. L'indice indique quel type est vivant, et le destructeur teste chaque type dans un fold. Si un constructeur lève une exception à l'intérieur de \texttt{emplace}, le variant ne contient rien.
\omcode{code/interviews-3/03-cpp/cpp/var.hpp}{8}{47}{Un petit variant.}\label{lst:iv3:cp:var}

Un pack dans \texttt{alignas} applique chaque alignement et retient le plus grand, de sorte qu'un \texttt{Wide} aligné sur 32 octets impose 32 à l'objet entier, ce que vérifie un \texttt{static\_assert} du test.

Le destructeur doit appeler le bon destructeur pour un type qu'on ne connaît qu'à l'exécution. Le fold se développe en un test par type~: \texttt{idx\_ == 0 ? destroy first : nothing}, puis \texttt{idx\_ == 1} et ainsi de suite, \texttt{i++} avançant le long du pack. \texttt{visit} fonctionne de la même façon. Il s'agit d'une chaîne linéaire de tests, qui convient pour quelques types. Les implémentations de \texttt{std::variant} génèrent en général une table ou un switch à la place.

\texttt{emplace} détruit d'abord l'ancienne valeur, pour que les deux ne partagent jamais le tampon, puis met \texttt{idx\_} à \texttt{none}. Si le constructeur lève ensuite une exception, \texttt{idx\_} reste à \texttt{none}, et le destructeur ne fait rien. C'est le même état que \texttt{valueless\_by\_exception}. Sans l'appel à \texttt{destroy}, l'ancien objet ne serait jamais détruit, et le compteur d'objets vivants du test le détecte.

\texttt{index\_of} est une lambda \texttt{constexpr} immédiatement invoquée, avec un fold sur le pack~: pour chaque type, s'il correspond, on enregistre \texttt{i}, puis on incrémente \texttt{i}. Un fold sur l'opérateur virgule s'évalue de gauche à droite, si bien qu'aucune récursion ni template auxiliaire n'est nécessaire. Si un type apparaît deux fois dans la liste, c'est la dernière correspondance qui l'emporte, alors que \texttt{std::variant} refuserait l'appel. Un type absent de la liste doit être refusé à la compilation, ce que fait le \texttt{static\_assert}. Sans lui, \texttt{emplace<std::string>} sur un \texttt{Var<int, char>} compile et déborde du tampon.

Le test exécute 30\,000 étapes aléatoires contre \texttt{std::variant}, avec un \texttt{monostate}, en couvrant un \texttt{emplace} qui lève une exception, et compare l'indice, le résultat de la visite et le nombre d'objets vivants, sous AddressSanitizer.

\iqwrong{«~Utilisez une \texttt{union} de \texttt{Ts.{}.{}.}.~» Une union ne peut pas être déclarée directement sur un pack. Elle demande un template récursif, ce qu'utilise la bibliothèque standard, et qui permet aussi \texttt{constexpr}.}
\end{solution}

\begin{solution}{iq:iv3:cpp:60}
\texttt{a + b} renvoie désormais un petit objet \texttt{Add} qui contient des références à \texttt{a} et \texttt{b}, et n'effectue aucun calcul. L'affectation à \texttt{d} lance une seule boucle, qui interroge l'arbre entier pour chaque élément, comme dans le \cref{lst:iv3:cp:expr}. Le test compte 0 allocation, là où la version simple en fait 2. \texttt{a = a + b} est sûr ici. \texttt{auto e} laisse une référence pendante, et AddressSanitizer signale \texttt{stack-use-after-scope}.
\omcode{code/interviews-3/03-cpp/cpp/expr.hpp}{6}{41}{Templates d'expression.}\label{lst:iv3:cp:expr}

Le type de \texttt{a + b + c} est \texttt{Add<Add<Vector, Vector>, Vector>}, donc toute l'expression se trouve dans le type, et le compilateur peut inliner \texttt{e[i]} jusqu'à \texttt{a[i] + b[i] + c[i]}. Le concept \texttt{Expr} empêche \texttt{operator+} de s'appliquer à des types quelconques.

\texttt{a = a + b} est sûr, car chaque élément ne dépend que du même indice, si bien que l'écriture de \texttt{a[i]} n'affecte jamais une lecture ultérieure. Une opération qui lit d'autres positions n'est pas sûre de cette façon~: un produit matriciel ou un décalage lit des éléments que la boucle a déjà écrits, raison pour laquelle ces bibliothèques font une copie dans ces cas.

Le piège tient à la durée de vie. \texttt{a + b + c} crée le \texttt{Add} intérieur comme temporaire, et l'extérieur en stocke une référence. Dans \texttt{d = a + b + c}, les temporaires vivent jusqu'à la fin de l'expression complète, boucle comprise. Avec \texttt{auto e}, le \texttt{Add} extérieur survit à la ligne, mais pas l'intérieur. Auparavant, \texttt{auto} sur un résultat \texttt{Vector} était anodin, ce qui explique qu'on écrive cela. Rien n'échoue à la compilation.

La réparation consiste à stocker un \texttt{Vector} nommé par référence, et une expression interne par valeur, puisqu'une expression n'est que deux références ou une copie imbriquée. Notre test choisit le type du membre avec \texttt{conditional\_t}, et la version avec \texttt{auto} s'exécute alors proprement sous AddressSanitizer, et donne 6. Les vecteurs nommés doivent encore survivre à \texttt{e}.

Les tailles incompatibles ne sont pas vérifiées ici, car \texttt{size} ne regarde que l'opérande de gauche. Une bibliothèque réelle les vérifie par une assertion. Nous n'avons mesuré aucun temps d'exécution.

\iqwrong{«~Le compilateur aurait de toute façon éliminé les temporaires.~» Les allocations de la version simple sont bien réelles, et le test en compte 2.}
\end{solution}
